% SOURCE docs/release/2026-10-29/source/docs/writing-repair/2026-10-29/src/b627a1d6f5e463b2.tex CHAR 0
\chapter{Interference, allocation policies and software adoption}
\label{v13:sw-networks}

The preceding relationship analysis showed why a policy may require joint
graph and action information. This chapter examines the allocation problem
itself. It distinguishes network and grouped targets, develops the two initial
software references, and compares spatial, experimental, learned-policy and
simulation alternatives. The recurring small graphs and cost denominators are
matched comparison devices; each implementation's different assumptions and
failure conditions remain part of the argument.

\section{The unit receiving treatment need not be the independent unit}

A household member's offer can influence another member's borrowing, and an
outreach campaign can affect connected firms. In such settings a potential
outcome depends on more than the individual's own action. Writing only
$Y_i(a_i)$ hides the policy being evaluated. A more explicit representation is
$Y_i(a_i,\mathbf a_{-i})$, often restricted through an exposure mapping of
other units' actions. The mapping and the network are scientific inputs;
neither is established merely by finding matching customer identifiers.

Two broad software constructions appear in the inspected sources. MossSpider
uses a network and exposure summaries in a network-targeted estimation
workflow. Inferference uses grouped formulas, allocation probabilities and
inverse-probability weighting. These constructions have different observation
laws and potentially different targets. Their package names do not establish
interchangeability. \citep{v13_sw_mossspider,v13_sw_inferference}

An allocation policy illustrates the target. Suppose two linked borrowers have
binary outreach actions $A_1,A_2$, independently assigned with probability $p$
under the proposed policy. Stipulate the synthetic financial outcome
\[
 Y_i(a_i,a_j)=100a_i+40a_j-30a_i a_j
\]
in CNY. The first term represents the own-action contribution, the second a
spillover and the last diminishing combined benefit. These are hypothetical
structural values, not bank estimates. Averaging under the allocation gives
\[
 V(p)=140p-30p^2.
\]
For $p=0.2$ the mean is 26.8 CNY; for $p=0.8$ it is 92.8 CNY. The allocation
contrast is 66 CNY per borrower. Exact enumeration of the four assignments
checks the formula. \citep{v13_project_exact}

The direct effect of changing only the individual's own action is instead
$100-30a_j$, or $100-30p$ when the neighbor is independently treated with
probability $p$. This differs from the total allocation contrast. Reporting
one as the other can misprice the campaign even when every arithmetic
calculation is correct. The bank must specify whether it changes one action,
the neighborhood's allocation or an entire branch policy, and which costs
are included in the outcome.

\section{MossSpider: network semantics are part of the estimand}

The inspected \texttt{NetworkTMLE} documentation describes expected incidence
under a treatment plan, with Monte Carlo integration over allocations for
stochastic plans. It states a direct-contact interference assumption. The
constructor rejects self-loops, noninteger node identities and nonbinary
exposure, creates a relabeled network while retaining original identities,
and constructs an unweighted adjacency matrix. \citep{v13_sw_mossspider}

These details matter for bank graphs. A weighted guarantee network is not the
same as an unweighted adjacency graph. Removing weights can turn a small
guarantee into the same exposure contribution as a large one. Relabeling nodes
is harmless only when results can be mapped back to the correct identities.
A linkage error can attach the correct numerical estimate to the wrong firm.
Semantic conformance therefore needs an identity roundtrip and an explicit
policy for unsupported weighted relationships.

Direction requires particular review. The documentation initially describes an
undirected graph but also explains a directed influence convention opposite
the graph arrow, and the selected constructor includes directed handling. This
combination requires a concrete test with an asymmetric small graph. One should
not infer direction from a diagram's appearance or assume that every library
uses source-to-target influence. A two-node one-way example with distinguishable
covariates can determine which neighbor contributes to each exposure summary.
Section~\ref{sec:exec24-interference-reuse} develops the selected mapping,
policy and variance paths. Their inspected direction convention still needs
matched native tests; directed bank-network adoption remains a proposed use. \citep{v13_sw_mossspider}

Degree restrictions also change the target. The documentation treats nodes
outside a requested degree range as background, with their exposures unchanged
by the intervention. This can be scientifically useful, but it is not the
same as deleting those nodes from the graph or intervening on all customers.
The response must state which nodes are in the target population, which remain
background and how their factual exposure affects the intervention outcome.

The source disclaims a missing-data procedure. Missing node attributes,
uncertain links and omitted customers are common in banking, but the existence
of an exposure-mapping function does not solve those observation problems.
One must either supply defensible complete data in the declared regime or
develop a separate missingness/linkage construction. Network-targeted estimation
does not automatically qualify ownership or inclusion truth.

The observed MIT text makes a restricted adapter worth considering. Full adoption
requires qualification of the exposure-model and mapping code, targeting step,
Monte Carlo integration and variance calculation, together with actual
dependency/runtime checks. The developed study retains deterministic-policy
refusals, background-node treatment and a potentially negative variance
quadratic; these restrictions belong in any proposed adapter. A claim of general contagion, dynamic financial-network equilibrium or
arbitrary spillovers would go beyond the selected implementation evidence.

\section{Inferference: grouped allocation and numerical integration}

The inspected \texttt{interference} function uses a formula that identifies
outcome, exposure, optional participation, propensity covariates and a grouping
variable. It accepts allocation probabilities and propensity-model choices
including \texttt{glm}, \texttt{glmer}, \texttt{HLfit} and an oracle route.
Its documented causal-estimation path is IPW with naive or robust variance
options. The selected code sorts records by group and checks group count,
allocation range and the variance option. \citep{v13_sw_inferference}

A group-based model can be appropriate when the scientific design justifies
within-group interference and the required between-group sampling assumptions.
It does not become a model of arbitrary connected borrowers by assigning
everyone a group label. Cross-group guarantees or shared market shocks can
violate the intended observation law. The grouping decision needs domain
evidence and should be subjected to sensitivity analysis where boundaries
are uncertain.

The default logit integrand is restricted to a single random effect in the
documented implementation. The code rejects more than one when that integrand
is selected. Thus a proposal with household and branch random effects needs a
different justified integration path; silently dropping one random effect is
a model change. An oracle propensity option is useful for synthetic reference
calculations, but supplied true parameters cannot be counted as evidence that
an estimated-propensity workflow is calibrated.

There is also a reproducibility issue visible directly in the source. The prose
describes \texttt{integrate\_allocation} as defaulting to false, while the
inspected function signature sets it to true. The proper response is to specify
and record the argument explicitly. Section~\ref{sec:exec24-interference-reuse}
traces the downstream ratio and derivative paths; their flag propagation,
allocation-mass and missing-group restrictions still require matched native
qualification. It is not appropriate to select whichever default gives a more
favorable result after comparing outputs. \citep{v13_sw_inferference}

For large groups, multiplying many probabilities can create numerical problems.
The documentation discusses an allocation-integration option in this context.
Qualification should include large and uneven groups, boundary allocation
probabilities, integration tolerances and variance derivatives. A successful
floating-point result does not prove adequate support; an underflow error does
not prove nonidentification. These are different failure modes and need
different diagnostics.

The observed DESCRIPTION declares GPL version 2 or later and a dated version.
That informs a dependency review but does not establish the license status of
every linked component or distribution arrangement. A candidate grouped-policy
reference should remain distinct from a decision to embed or port it.

\section{A common comparison must distinguish network and group assumptions}

The two packages can be compared on a common small allocation problem only
after their target, exposure mapping and sampling assumptions are aligned.
For the two-borrower example, the exact four-assignment enumeration supplies a
transparent functional reference. A larger benchmark should vary degree,
group size, allocation probability, support, observation errors and dependence.
It should include examples where a direct-contact or group restriction is false.

Metrics should separate the target mean, direct and spillover contrasts,
coverage, finite availability, Monte Carlo integration error and total cost.
Increasing simulated allocations reduces integration error conditional on a
fitted model. It does not create independent groups or validate the graph.
Uncertainty must follow the actual sampling unit; counting nodes as independent
merely because the graph is stored in a table can produce misleading precision.

The implementation contract should retain the original graph, its direction
and weights, the identity map, intervention population, allocation rule,
background nodes and exposure mapping. A result should be rejected as
semantically incompatible if a proposed bridge silently changes any of these.
This requirement is shared by library calls, command-line jobs and agent tools.
An interface can make a scientific restriction visible; it cannot make it true.

\section{Provisional reuse choices and the evidence needed to change them}

The inspected software supports a portfolio of candidate components rather
than a single universal engine. The following choices describe potential uses and qualification conditions,
not accepted integrations or benchmarks. The earlier functional and
longitudinal chapters develop the corresponding source findings
(Chapters~\ref{v13:sw-functional} and~\ref{v13:sw-regimes}); the table retains
the conditions that must survive implementation rather than treating those
studies as still absent.

\begin{longtable}{P{0.18\textwidth}P{0.34\textwidth}P{0.39\textwidth}}
\toprule
Component & Candidate role & Conditions to preserve and qualify\\
\midrule
\endhead
DirectLiNGAM & Functional discovery adapter & Reviewed base-class behavior, prior orientation,
functional assumptions and adverse latent/Gaussian cases\\
y0 & Symbolic ID/IDC adapter & Reviewed recursive paths, conditional-path restriction,
graph validation, expression roundtrip and nonidentification references\\
delicatessen & Equation and numerical reference & Solver, Jacobian, covariance,
sampling unit and complete dependency contract\\
statsmodels & Point-treatment GMM baseline & Matched population, propensity bounds,
actual covariance behavior and dependent-data limits\\
pygformula & Longitudinal-regime adapter & History/event meaning, simulation,
bootstrap unit, diagnostics and policy scope\\
grf & R/native heterogeneous-effect component & Native algorithms, variance,
aggregation measure, ABI and license/distribution review\\
CausalGPS & Continuous-density component & Density/transformation helpers,
counterfactual doses, support and downstream estimation\\
causaldrf & Dose-response reference & Full function, interval procedures,
support, recovered holder notice and dependency terms\\
MossSpider & Restricted network-policy component & Direction, mapping, identity,
targeting, variance and observation restrictions\\
inferference & Grouped-allocation reference & Integrands, numerical defaults,
variance, group law and dependency terms\\
\bottomrule
\end{longtable}

An integration decision is stronger than identifying a promising function.
Direct adoption needs a stable, compatible dependency and scientific contract.
An adapter needs semantic parity and error preservation. A port needs the right
to reuse the implementation, an upstream behavioral oracle, independent
references, numerical/statistical equivalence criteria and a maintenance owner.
A research reference can be valuable without becoming a product dependency.
Deferral is justified when these prerequisites are missing; rejection should
state a concrete incompatibility and a revisit condition.

Commercial comparison must apply the same discipline. Vendor documentation can
establish a dated documented feature, but not measured agreement or operating
cost without an executed comparison. An inaccessible manual cannot support a
claim that the vendor lacks an operation. Our proposed graph-to-effect or
banking workflow is likewise not a delivered competitive advantage. A sound
build/buy/partner decision compares the vendor's documented and measured state
with our implemented, conditional, proposed or absent state, and makes the
missing evidence explicit.

Licensing is one part of this decision. The selected MIT and BSD texts,
repository GPL text and package license declarations provide traceable starting
points. Dependencies, bundled data, separate notices, native libraries and the
intended distribution arrangement still need review. A subprocess boundary
does not automatically settle licensing obligations, just as a Python wrapper
does not establish that a native operation has been qualified.

The software survey can therefore guide implementation immediately without
claiming premature adoption. It identifies precise operations, preserves their
mathematical meaning and specifies the evidence needed for the next decision.
Runtime, bank validation and independent scientific acceptance remain separate
from these source-grounded provisional choices.


% SOURCE docs/release/2026-10-29/source/docs/writing-repair/2026-10-29/src/89ab4d3132a15fc3.tex CHAR 0
\section{Interference software must preserve the population and the policy}
\label{sec:exec24-interference-reuse}

A bank considering outreach to connected borrowers must decide whose outcome it
wants to improve and which assignments it can change. A household contact, a
branch allocation and a directed financial relationship define different
interventions. MossSpider 0.0.3 and inferference 1.0.3 provide useful, but
different, implementation references. The former exposes a network-based
targeted procedure; the latter constructs grouped allocation inverse-probability
weights. Their selected commits are respectively
\texttt{62e6bd20364c7ef01b6964fd79f21be84e5ac579} and
\texttt{b0c785e3f60541b993ff4e7324672b761a5f9723}. The comparison below is
an inspection of authentic source, not a native benchmark or an assertion that
their unrestricted statistical ambitions have been qualified.
\citep{exec24_interference_mossspider,exec24_interference_inferference}

\subsection{A network population is not automatically a population of groups}
Let $A_i$ be a binary contact decision and let $G_{ij}=1$ mean that borrower
$j$'s treatment contributes to the exposure summary for borrower $i$. Define
$S_i(A)=\sum_jG_{ij}A_j$. A proposed network policy specifies a law for the
whole assignment vector and a population $E$ of eligible outcomes. A candidate
target is
\[
 \psi_E(\pi)=\frac{1}{|E|}\sum_{i\in E}E_\pi[Y_i(A_i,S_i(A))].
\]
This notation states a target; it does not identify it. Consistency, the
adequacy of the exposure mapping, measured confounding, support and a justified
sampling or network law remain substantive requirements. An API accepting an
adjacency matrix supplies none of them.

The selected MossSpider constructor requires integer node identifiers, refuses
self-loops and nonbinary treatment, saves an original-identifier mapping and
constructs an unweighted adjacency matrix. Its exposure helper multiplies rows
of that matrix by the ordered treatment vector. For the directed chain
$0\to1\to2$ and treatment vector $(1,0,0)$, row-neighbor exposure is
$(0,0,0)$; using incoming neighbors gives $(0,1,0)$. Both are calculable, but
they answer different banking questions. An adapter must declare the meaning
of an arrow, preserve node ordering and test an asymmetric graph. Automatically
symmetrizing a relationship network changes the intervention.
\citep{exec24_interference_mossspider,exec24_interference_reference}

For grouped partial interference, let group $g$ contain $n_g$ people and let
$\mu_g(\alpha)$ be its average outcome under allocation $\alpha$. Two
reasonable targets are
\[
 \psi_{\mathrm{group}}=\frac1M\sum_{g=1}^M\mu_g(\alpha),
 \qquad
 \psi_{\mathrm{person}}=\frac{\sum_g n_g\mu_g(\alpha)}{\sum_g n_g}.
\]
The inspected inferference point-estimate path first constructs group means
and then averages their weighted contributions with denominator $M$. Consider
two groups of sizes two and eight with stipulated outcome means USD~100 and
USD~0. Equal-group averaging gives USD~50, whereas equal-person averaging gives
USD~20. Neither is a software error. Substituting one for the other would be a
product error if a bank's objective were total borrower benefit. The treatment-
specific group contribution also divides by total group size, rather than the
number treated; its later allocation normalization must remain part of the
same estimand. An integration cannot compare the two packages simply because
both return a field called an effect.
\citep{exec24_interference_inferference,exec24_interference_reference}

\subsection{Eligibility restrictions change the intervention itself}
MossSpider's degree restriction selects the outcome population and also affects
policy simulation. The inspected generator keeps excluded nodes at their
observed treatment while resampling eligible nodes. This is a mixed policy,
not an all-network stochastic intervention. In a star network, suppose the hub
is excluded by a degree threshold, its observed treatment is one, and each
eligible leaf has stipulated outcome $10A_{\mathrm{hub}}$. Keeping that hub
fixed gives ten per leaf. Resampling the hub with probability $1/4$ gives
$5/2$. The distinction persists even if only the leaves' outcomes are averaged.
Eligibility, the treatment of excluded neighbors and the resulting target
population must therefore travel together in the request and response.
\citep{exec24_interference_mossspider,exec24_interference_reference}

Boundary behavior also requires an explicit contract. The selected fit
documentation describes a one-sample shortcut for deterministic policies, but
the selected body rejects uniform all-zero and all-one plans. It rejects an
integer argument even when its mathematical value could be a probability. A
wrapper should expose these selected-version restrictions, rather than promise
the documented shortcut or silently perturb zero into a small positive number.
The source includes tests for several input and model refusals; this inspection
does not execute those tests or claim that every boundary is covered.
\citep{exec24_interference_mossspider}

\subsection{Allocation mass must appear exactly once}
For a group assignment $a$, denote its observed probability by $p_\theta(a)$
and its proposed policy probability by $p_\alpha(a)$. The desired importance
ratio is $w_\theta(a;\alpha)=p_\alpha(a)/p_\theta(a)$ when both probabilities
refer to the same assignment and conditioning information. The inspected
default inferference integrand already divides the observed Bernoulli factors
by the policy factors. Consequently its integrated value represents a
probability ratio under that default construction. The reciprocal branch and
the branch that multiplies by policy mass are not interchangeable.
\citep{exec24_interference_inferference}

An exact no-random-effect witness makes the normalization visible. For two
treated members with common observed probability $q=1/2$ and proposed
probability $\alpha=1/4$, observed mass is $1/4$ and policy mass is $1/16$.
The ratio integral is four, so its reciprocal is $1/4$. Multiplying by policy
mass again gives $1/64$. If $q$ is parameterized by a common logit coefficient
$\beta$, then
\[
 \frac{\partial w}{\partial\beta}=-2(1-q)w,
\]
which equals $-1/4$ in the first construction and $-1/64$ in the duplicated-mass
construction. This is an authored algebraic example of two different
normalizations, not an observation of an executed package failure.
\citep{exec24_interference_reference}

The implementation review must follow the entire argument chain. The selected
derivative-array helper passes a false allocation-integration flag to its
derivative helper. That helper consumes the formal argument but does not place
it in the constructed gradient arguments; the downstream weight function then
uses its default. Reading only the first call would therefore misdescribe the
executed semantics. In a future native qualification, both flag values, the
default and a custom integrand must be checked against a known integral and an
independently calculated derivative. The selected repository includes numerical
weight and derivative tests, but their existence alone does not settle every
flag/default combination. A product should record the resolved settings and
refuse a score/weight mismatch.
\citep{exec24_interference_inferference}

\subsection{Numerical output does not settle uncertainty or missingness}
The grouped weight calculation catches an integration error and returns a
missing value. Later group aggregation uses a sum with missing values removed
while retaining the original group count. If three stipulated group
contributions are $10$, missing and $30$, this expression yields $40/3$. The
complete-case mean is $20$, and a hypothetical complete contribution of $60$
for the missing group would give $100/3$. None is a justified replacement for
the missing scientific contribution without further assumptions. The wrapper
must report submitted and usable groups, retain the failure and qualify or
refuse the target; it must not reinterpret a finite result as complete data.
\citep{exec24_interference_inferference,exec24_interference_reference}

The inspected grouped robust calculation also solves a score crossproduct
matrix, and the returned intervals use normal quantiles. A matrix inverse
requires numerical diagnostics; a label such as robust requires a justified
group-sampling law. In the network implementation, one variance uses squared
weighted residuals and another uses their adjacency quadratic form with the
diagonal filled. A symmetric path with residual vector $(1,-2,1)$ gives
residual-square expression $2/3$ and adjacency expression $-2/9$. This shows
that the bare finite-sample quadratic need not be positive semidefinite. It
does not show that these residuals arose from the fitted procedure or refute
any stated population theorem. Negative or undefined variance is nevertheless
a necessary refusal/diagnostic fixture, alongside simulation under a separately
justified network law. Clipping it to zero would conceal the issue.
\citep{exec24_interference_mossspider,exec24_interference_inferference,exec24_interference_reference}

\subsection{Separate reuse decisions and a bounded public prototype}
The proposal retains both packages as references. An original-language worker
is a conditional prototype route: it avoids an immediate scientific rewrite,
but still needs exact identity, default, missingness, policy and uncertainty
conformance. An exact finite-policy oracle supplies a useful baseline. A
whole-library port is not justified by the present evidence. It would add
maintenance and scientific obligations before demonstrating an advantage.
The two reference choices do not approve either as a general contagion or
dynamic banking network engine.
\citep{exec24_interference_reference}

MossSpider's inspected MIT notice requires notice preservation and disclaims
warranty; its selected setup lists several dependencies without fixed versions.
inferference's DESCRIPTION declares GPL version two or later and identifies its
R dependencies and maintainer. These are source facts, not a complete
distribution-rights determination. A platform owner must review the actual
installed distribution and transitive dependencies before shipping. Selected
commit and publication dates establish historical identity, not future support.
\citep{exec24_interference_mossspider,exec24_interference_inferference}

The first public prototype should prespecify asymmetric direction, unequal
group size, excluded-hub, deterministic-boundary, known-integral, derivative,
missing-group and singular/negative-variance fixtures. Scientific owners must
define the target and sampling law; adapter owners preserve identities and
failure states; evaluation owners separate native conformance from uncertainty
coverage. Every attempt contributes to cost. For a stipulated USD~2,400 total,
100 submitted requests and 60 usable outputs imply USD~24 per submission and
USD~40 per usable output. These figures illustrate denominators, not vendor
prices or bank returns. A later expensive bank experiment needs a qualified
public prototype and a credible economic case; neither follows merely from
the availability of these two interfaces.
\citep{exec24_interference_reference}


% SOURCE docs/release/2026-10-29/source/docs/writing-repair/2026-10-29/src/82b72f77cea3839f.tex CHAR 0
\section{Commercial network tools: spatial impacts, exposure designs and switchbacks}
\label{exec20:network-scope}

A branch intervention can affect customers outside the treated branch, while a
borrower's funding choice can propagate through guarantees and counterparties.
Commercial software already offers several relevant routes, but they answer
different questions. A spatial regression estimates a response under a supplied
weight matrix and model; an exposure experiment compares outcomes under an
assignment design; a switchback changes a population's treatment over time.
The proposal must compare these routes at the level of their targets, information
and assumptions. A graph interface or a feature called ``impact'' is not a
common causal estimand.

The selected primary documentation supports concrete profiles for Stata,
SAS/ETS, ArcGIS Pro and Statsig. Their strengths include supplied-network
estimation, spatial data preparation, diagnostic output, experimental assignment
and analysis services that our proposed library has not implemented. The
comparison therefore asks which original engines or services merit a public
prototype, and which semantics the proposed common interface must preserve.
It does not claim existing product parity.
\citep{exec20_network_stata_spregress,exec20_network_sas_latest_spatial,
exec20_network_arcgis_tool,exec20_network_statsig_types}

\subsection{Define the intervention before comparing an impact}
\label{exec20:network-target}

Let $Y$ be an $n$-vector of branch outcomes, $a$ a vector of intervention levels,
and $W$ a supplied matrix whose entry $W_{ij}$ weights the influence of unit $j$
on unit $i$. This receiver-row convention differs from the source-row adjacency
used by some discovery packages. For a stipulated linear structural equilibrium,
consider
\begin{equation}
 Y=\rho WY+\beta a+\gamma Wa+U,
 \qquad
 S(W)=\frac{\partial E[Y\mid\operatorname{do}(a)]}{\partial a^{\mathsf T}}
      =(I-\rho W)^{-1}(\beta I+\gamma W).
 \label{exec20:network-response}
\end{equation}
The second expression requires an invariant structural mechanism, an exogenous
intervention in $a$, a fixed relevant $W$, and an invertible feedback operator.
The sufficient stability restriction $|\rho|r(W)<1$, where $r(W)$ is the
spectral radius, also gives the convergent propagation expansion
$S=\sum_{k\geq0}\rho^kW^k(\beta I+\gamma W)$. A fitted conditional spatial
mean can have the same algebra without supplying those intervention premises.
An endogenous treatment, omitted common cause or treatment-dependent network
requires additional identification information. The equation is a project
derivation for the stated model, not a theorem about every vendor estimator.
\citep{exec20_network_reference}

The average own derivative is $D=\operatorname{tr}(S)/n$, the average total
derivative is $T=\boldsymbol{1}^{\mathsf T}S\boldsymbol{1}/n$, and the
average spillover derivative is $T-D$. These average quantities also differ from
the response to treating one selected branch $j$, which is the column $S e_j$.
Stata's \texttt{estat impact} documents the averages of own, off-diagonal and
total derivatives of its reduced-form mean, and exposes their Jacobians and
covariance matrices. ArcGIS Pro 3.5 instead documents ``simple impacts,'' with
the direct impact equal to the coefficient. A common result named
\texttt{direct\_effect} must preserve which definition was used.
\citep{exec20_network_stata_impacts,exec20_network_arcgis_method}

For an exact banking reference, set $\gamma=0$, $\beta=10$ dollars per
intervention unit, $\rho=1/2$, and
\begin{equation}
 W=\begin{pmatrix}0&1\\1&0\end{pmatrix},\qquad
 S=\begin{pmatrix}40/3&20/3\\20/3&40/3\end{pmatrix},\qquad
 (D,T-D,T)=(40/3,20/3,20).
 \label{exec20:network-feedback}
\end{equation}
At an assumed cost of USD 12 to intervene at branch 1, the raw coefficient
gives a margin of USD $-2$, the own-branch response gives USD $4/3$, and the
two-branch portfolio response gives USD 8. These are three target calculations,
not contradictory estimates of one quantity. If the actual tie from branch 1
to branch 2 is missing from the supplied matrix, the modeled portfolio response
becomes 10 and its margin becomes USD $-2$. Neither more optimizer iterations
nor a generic robust covariance option restores the omitted path. The example
uses an authored structural law and assumed costs; it is not observed bank or
commercial-runtime evidence. \citep{exec20_network_reference}

\subsection{Weights, identities and uncertainty are substantive choices}
\label{exec20:network-weights}

Stata distinguishes scalar normalization from row normalization. Dividing $W$
by a common scalar and changing spatial coefficients correspondingly can
preserve the model. Normalizing each row separately generally changes its
meaning. The manual also refuses a sample smaller than the matrix's unit set
unless \texttt{force} is requested, warning that removing units changes spatial
connections and normalization. Its maximum-likelihood robust covariance is
described for nonnormal independent and identically distributed innovations;
the heteroskedastic GS2SLS option relaxes identical distribution while retaining
independence. None of these choices supplies arbitrary residual network
dependence. \citep{exec20_network_stata_spregress}

The exact three-branch star
\begin{equation}
 W_\star=\begin{pmatrix}0&1&1\\1&0&0\\1&0&0\end{pmatrix},\qquad
 W_{\rm row}=\begin{pmatrix}0&1/2&1/2\\1&0&0\\1&0&0\end{pmatrix}
 \label{exec20:network-normalization}
\end{equation}
illustrates the distinction. With $\rho=1/4$, $\beta=10$ and $\gamma=0$, a unit
intervention at branch 1 has responses $(80/7,20/7,20/7)$ under $W_\star$ and
$(32/3,8/3,8/3)$ under $W_{\rm row}$. In contrast, replacing $W_\star$ by
$W_\star/2$ and $\rho$ by $1/2$ preserves the original response exactly.
The two row conventions describe different propagation assumptions. An agent
must not select one merely to match a package default.
\citep{exec20_network_reference}

SAS/ETS's selected SPATIALREG manual, labeled 2023.10 and dated August 2024,
documents a broad cross-sectional model family including SAR, spatial Durbin,
error and moving-average specifications. Its default row standardization
replaces missing weights with zeros; \texttt{NONORMALIZE} instead refuses missing
off-diagonal entries and treats diagonal entries as zero. A missing relationship
is not scientifically equivalent to a known absent relationship. The
\texttt{SPATIALID} field aligns observations with the matrix, but numeric
noninteger IDs are truncated for matching. The proposed interface should use
validated immutable identifiers and an explicit missing-tie policy before
dispatch, preserving all exclusions and conversions.
\citep{exec20_network_sas_latest_spatial}

The selected SAS SAR likelihood uses independent normal innovations. It offers
Hessian, outer-product and QML covariance choices, with Hessian as default and
additional restrictions when certain likelihood approximations are selected.
Stata's default impact delta method conditions on regressors; its unconditional
option includes their sampling variation and has sample-selection restrictions.
These documented distinctions are useful requirements for an adapter, but
do not by themselves prove identification or useful finite-sample coverage.
\citep{exec20_network_sas_latest_spatial,exec20_network_stata_impacts}

Full parameter covariance must also survive extraction. In the two-branch
reference, $T=\beta/(1-\rho)$ has gradient $(2,40)$ in the order $(\beta,\rho)$.
For the illustrative covariance
\begin{equation}
 V=\begin{pmatrix}1&1/20\\1/20&1/100\end{pmatrix},\qquad
 \nabla T^{\mathsf T}V\nabla T=28,
 \label{exec20:network-variance}
\end{equation}
whereas dropping the cross-covariance gives 20. The covariance is assumed and
positive definite; the calculation is an exact delta-method reference, not an
estimated standard error or coverage finding. It explains why a CLI or MCP
response containing only coefficients and individual standard errors cannot
generally support the intended propagated-value calculation.
\citep{exec20_network_reference}

Stata's spatial-panel route further requires strongly balanced panels. Its
fixed-effects formulation eliminates panel effects and cannot fit covariates
constant within a panel; random-effects formulations impose their own
distribution and independence restrictions. The displayed innovations are
independent across panels and time. This is a concrete panel capability worth
comparing, but it does not automatically qualify evolving guarantee networks or
dependent servicing histories. \citep{exec20_network_stata_panel}

ArcGIS Pro adds geographical neighborhood construction, diagnostic maps,
lag/error/combined model choices and \texttt{arcpy.stats.SAR}. Its documented
defaults and inferred bandwidths need to be retained with the result. The
method page describes robust errors and an AK-test-triggered spatial HAC option,
using a separate nearest-neighbor triangular-kernel matrix. It also documents a
30-percent connectivity refusal. That is a rule of the selected tool, not a
universal boundary for causal validity. Basic, Standard and Advanced availability
is listed, but the listing does not establish redistribution, embedding or
deployment rights. The selected prose is inspected; image-only equations and
flowcharts still need a separate original-source review.
\citep{exec20_network_arcgis_method,exec20_network_arcgis_tool}

\subsection{Exposure support is a property of the experiment}
\label{exec20:network-exposure}

For a different route, let $Z_i$ be treatment assignment and let
$E_i(Z,G)$ summarize the assignments of unit $i$'s neighbors in a declared
network $G$. An exposure response $Y_i(z,e)$ assumes that this summary retains
the relevant treatment interference. Even if the graph is measured perfectly,
the assignment design may never generate some required $(z,e)$ combinations.
For two connected borrowers with exactly one treated, the only possible
exposures are $(1,0)$ and $(0,1)$, each with probability $1/2$ for a given
borrower. The cells $(0,0)$ and $(1,1)$ have zero probability. Direct effects at
zero neighbor treatment cannot be recovered by weighting observations into a
cell the experiment never visits. \citep{exec20_network_reference}

Independent half-probability treatment assignments instead give each exposure
cell probability $1/4$. With known positive probabilities
$\pi_i(z,e)=\Pr\{Z_i=z,E_i=e\}$, an authored finite-population mean reference is
\begin{equation}
 \widehat\mu(z,e)=\frac1n\sum_{i=1}^n
 \frac{\mathbf{1}\{Z_i=z,E_i=e\}Y_i}{\pi_i(z,e)}.
 \label{exec20:network-ht}
\end{equation}
Under consistency for the declared exposure mapping and the known assignment
law, linearity of expectation yields the mean potential outcome for that cell.
Enumerating all four assignments for the stipulated law $Y_i(z,e)=5+2z+3e$
gives expected means 5, 8, 7 and 10 for $(0,0),(0,1),(1,0),(1,1)$ respectively.
The direct contrast at $e=0$ is 2, and the spillover contrast for untreated units
is 3. This verifies a point-mean identity. It supplies neither a useful variance
estimate from two borrowers nor a general confidence interval; joint exposure
probabilities, dependence and the sampling/design law remain necessary.
\citep{exec20_network_reference}

Statsig's selected cluster-unit documentation gives a different, concrete
service: assign at organization level and analyze user revenue using a specified
one-to-many organization/user relationship. Its ratio treatment accounts for
numerator-denominator covariance with the delta method. That can address the
distinction between assignment and analysis units, but does not establish that
all interactions lie within organizations, that treatment cannot alter the
denominator, or that the desired exposure cells have positive probability.
Those questions must remain in the proposed experiment contract.
\citep{exec20_network_statsig_units}

\subsection{Switchbacks change interference and time together}
\label{exec20:network-switchback}

Statsig's legacy technical documentation describes switchbacks in which
populations or regions switch treatment by time window. It documents event
attribution using exposure timestamps, unit IDs and an attribution window;
bucket-level and variant-level metrics; separate resampling of treatment and
control buckets 10,000 times; and percentile confidence intervals. Configuration
includes schedules, window lengths and burn-in/burn-out exclusions, with
randomized bucketing described as an advanced feature. These are substantive
capabilities for a public experimental comparison, not merely a network label.
\citep{exec20_network_statsig_types,exec20_network_statsig_switchback}

Changing everyone in a region together can reduce one form of simultaneous
cross-treatment interference while introducing an important time comparison.
For an authored law $Y_t(z)=100t+2z$, observing control at $t=0$ and treatment at
$t=1$ gives an observed difference of 102 although the immediate causal effect
is 2. With carryover $Y_t=2Z_t+4Z_{t-1}$ under an alternating schedule, treated
windows follow control and yield 2, while control windows follow treatment and
yield 4. The observed difference is $-2$, the immediate effect is 2, and the
sustained treatment effect is 6. Resampling the same contaminated windows cannot
create missing assignment or washout information. These are explicit
counterexamples to an unrestricted interpretation, not claims of an observed
Statsig failure. \citep{exec20_network_reference}

The prototype should therefore declare the policy target, randomization of
start phases and schedules, carryover horizon, attribution overlap and effective
resampling unit. Discarding exposures during burn periods is not by itself a
demonstration that all associated outcomes are free of earlier treatment.
Adjacent buckets or shared regional shocks require a justified dependence
treatment. The saved September 2025 documentation also lists profile limits,
including unavailable time-series and dimension breakdowns and unavailable CUPED
and sequential testing for that legacy switchback profile. Those statements
must not be promoted to claims about current commercial availability.
\citep{exec20_network_statsig_switchback}

\subsection{A comparison that can guide the first prototype}
\label{exec20:network-comparison}

For an initial equilibrium-model prototype, Stata and SAS provide documented
original engines with different model, covariance and weight policies. ArcGIS
offers a geographical workflow and a Python surface, with an impact convention
that needs explicit reconciliation. Statsig supplies an experimental alternative
when population/time assignment is the intended intervention. A rational first
comparison chooses a declared target and then asks which of these profiles
actually supports it; it does not rank unlike outputs by one accuracy number.
Our present implemented contribution is the exact authored references. Model
fitting, licensed execution, native covariance behavior and service integration
remain proposed capabilities.

The public prototype can begin with the two- and three-branch laws above,
generated assignments, missing-tie perturbations and explicit carryover. It
should compare target preservation, refusals, graph and sample identities,
uncertainty definitions and economic decisions before scale. Public empirical
datasets can test import, geography and workflow behavior, but their observed
network alone supplies no intervention truth. Actual runtime, memory, analyst
effort, licensing, data movement and upgrade costs must be measured rather than
inferred from documentation. Commercial adapters should preserve the original
engine where justified; porting a proprietary algorithm is neither necessary
nor supported by the present evidence.

Coverage remains incomplete for the exact editions and interference targets of
Optimizely, EViews, IBM SPSS and the commercial probabilistic-network families.
The requested Optimizely route returned a general homepage, which supports no
technical capability claim. Supplied Bayesian networks and ordinary pooled
regression must not stand in for an uninspected exposure estimator. Their next
steps are concrete: identify the exact procedure and version, acquire the
original technical specification and rights terms, and compare its target with
the same public reference. None is classified as lacking a feature merely
because its original has not yet been reviewed. These unresolved profiles and
unknown prices limit the market comparison while allowing the supported
prototype designs to proceed.


% SOURCE docs/release/2026-10-29/source/docs/writing-repair/2026-10-29/src/ff4336dd688ce96a.tex CHAR 0
\section{R network workflows: the exposure, assignment and population must travel together}
\label{sec:exec25-r-network}

A bank choosing which branches receive a retention intervention faces several different network questions. It may ask for an exposure-specific response under a randomized allocation, an effect among customers whose treatment changes with an encouragement, a test of a sharp exposure null, or the value of a policy selected from a constrained class. These are different targets. A common software interface must preserve their distinctions instead of treating every object called a spillover effect as an interchangeable number. The network itself is part of the observation contract: borrower identifiers, direction, weights, missing relationships and the time at which an edge became available determine which exposure is being measured.

The present comparison extends the earlier inferference review with eight R implementations at authentic pre-cutoff revisions. The selected versions are interference 0.1.0, latenetwork 1.0.1, BicliqueRT 1.0.0, NetworkTargeting 0.0.0.9000, NetworkCausalTree 1.0.0, MisspecifiedInterference 1.2, autognet 0.0.3 and clusteredinterference 1.0.1. Their inspected revisions range from 2019 through June 2026 and precede the client cutoff of 7 October 2026. Repository dates and declarations identify source bytes; none of these selected packages or their tests was executed in this comparison. The findings below concern inspected expressions, composed call paths and project-authored exact witnesses. They are not installed-package benchmarks or imported proofs from papers named in source comments.~\citep{exec24_interference_inferference,exec25_rnet_exposure,exec25_rnet_late,exec25_rnet_biclique,exec25_rnet_target,exec25_rnet_tree,exec25_rnet_nmr,exec25_rnet_auto,exec25_rnet_cluster}

\subsection{Design-based exposures: a label is a mathematical contract}

Let $Z=(Z_1,\ldots,Z_n)$ denote the assignment vector, $A$ the declared adjacency matrix, and $e_i(Z,A)$ a specified exposure map for customer $i$. A mean under exposure $c$ requires an identifying relationship between the potential response and that map, as well as support under the assignment law. The notation $\pi_i(c)=\Pr\{e_i(Z,A)=c\}$ concerns that law; it is not an ordinary propensity inferred merely from the presence of a network field. Joint probabilities $\pi_{ij}(c,d)$ also matter for uncertainty. A command must retain the design or simulator, the map, and the precise population along with the resulting estimates.

The interference package constructs first-degree categories from own assignment and a positive peer-exposure indicator. Its selected expression multiplies the assignment row vector by the adjacency matrix. Thus an asymmetric adjacency convention must be checked against the bank's arrow convention. Its two-hop option uses a thresholded matrix square with a zeroed diagonal. A two-step path is not necessarily a neighbor at shortest-path distance exactly two: triangles can create overlap with direct neighbors. A translation layer must describe that convention rather than silently replace it with a different graph operator.~\citep{exec25_rnet_exposure}

One inspected full-neighborhood helper is particularly useful as a qualification target. Its documentation names two conditions: the focal unit and all neighbors treated, or the focal unit and all neighbors controlled. The selected implementation instead defines the second indicator by exposure count less than neighborhood size. In a two-node connected graph with assignment $(1,0)$, neither closed neighborhood is all controlled, yet both are not all treated. The authored witness therefore produces two true selected control indicators and two false literal all-control indicators. The distinction is substantive: a comparison with a mixture of partially treated neighborhoods answers a different question from comparison with universal local control. Admission should wait for an explicit choice of contract and a version-specific native test.~\citep{exec25_rnet_exposure,exec25_rnet_choice}

The probability builder also makes an explicit approximation. With $R$ simulated assignments and exposure-indicator matrix $I_c$, the same-exposure matrix is
\[
 \widehat P_{cc}=\frac{I_cI_c'+I_n}{R+1},
 \qquad
 \widehat P_{cd}=\frac{I_cI_d'}{R}\quad(c\ne d).
\]
The additive diagonal correction is not the exact enumerated design probability. If four assignments exhaust a two-person Bernoulli design and each of four categories occurs once for a focal person, the four corrected marginals are $2/5$ and sum to $8/5$, whereas the exact probabilities are $1/4$. This does not by itself invalidate every use of a deliberate Monte Carlo correction. It requires the adapter to distinguish corrected finite-simulation quantities from an exact probability distribution and to assess their effect on estimates and uncertainty. A positive correction for an unseen event does not establish that the scientific design supports that event.~\citep{exec25_rnet_exposure,exec25_rnet_choice}

The selected permutation generator rounds $np$ to a fixed treated count, samples assignments and rejects duplicates unless repetition is allowed. That is a fixed-count design, not independent Bernoulli assignment. The estimator then distinguishes a Horvitz--Thompson total from a H\'ajek normalized total, and marks an exposure with no observed contributors as unavailable. The proposed reuse is a restricted design-based adapter with explicit mapping and probability modes. Direct adoption of a generic exposure label without these choices is rejected. Complete covariance helpers, random-state behavior and dependencies remain qualification work.~\citep{exec25_rnet_exposure,exec25_rnet_choice}

\subsection{Instrumented network responses: support and the variance law remain separate}

The latenetwork direct path forms empirical joint strata of an instrument $Z_i$ and an instrumental exposure $T_i$. It compares conditional outcome and treatment means across instrument values, then returns their ratio. Schematically the ratio is
\[
 \widehat\theta(t)=
 \frac{\widehat E[Y\mid Z=1,T=t]-\widehat E[Y\mid Z=0,T=t]}
      {\widehat E[D\mid Z=1,T=t]-\widehat E[D\mid Z=0,T=t]}.
\]
For a bank encouragement, assignment validity, exclusion, treatment-response restrictions and the interpretation of the exposure stratum are substantive requirements. An empty empirical stratum or near-zero denominator is a separate numerical and inferential problem. Neither is resolved by a function returning a local-effect label.~\citep{exec25_rnet_late}

The selected validator rejects missing inputs, nonbinary assignments, asymmetric adjacency and nonzero diagonal entries. In particular, its symmetry test counts unequal transpose entries; it does not sum signed differences that would cancel. The direct routine constructs undirected graph distances on the original population before restricting to the requested subpopulation. Paths through customers outside that subpopulation can therefore remain in the distance matrix. This differs from computing distances on the induced subgraph and must be carried into the uncertainty contract.~\citep{exec25_rnet_late}

The inspected HAC branch uses $K_{ij}=\mathbf 1\{d(i,j)\le b\}$ in a score quadratic, followed by a square root. A distance cutoff matrix need not be positive semidefinite. For the three-node path at bandwidth one and a centered illustrative score,
\[
 K=\begin{pmatrix}1&1&0\\1&1&1\\0&1&1\end{pmatrix},
 \qquad v=(1,-2,1)',\qquad \frac{v'Kv}{9}=-\frac{2}{9}.
\]
This exact matrix witness shows a property of the selected quadratic, not a claim that a particular native fitted sample produces this score. The wild-bootstrap path uses a different neighborhood-overlap covariance and a separate small-eigenvalue rule. These are not interchangeable uncertainty services. Qualification must fix graph distances, bandwidth, score law, weak-instrument handling and all nonfinite results before comparing interval coverage. The proposal retains latenetwork as a bounded reference and adapter candidate, with no universal network-IV or dependence guarantee.~\citep{exec25_rnet_late,exec25_rnet_choice}

\subsection{Conditional randomization tests need typed failure and tie semantics}

BicliqueRT accepts a design function, an exposure function and a null-equivalence relation. The inspected preparation places the observed assignment in the first column, simulates other assignments from the supplied design, and builds a null-exposure graph. The selected test computes a statistic on the focal units and assignments of the chosen biclique. The native bimax route has its own compiled call and bounded result capacity. These are useful building blocks for a declared conditional test; the package title does not establish validity for an arbitrary assignment sampler or an incorrectly specified null.~\citep{exec25_rnet_biclique}

The return contract needs special attention. Missing statistics are returned with a numerical p-value field equal to three; a degenerate randomization distribution can return two. These values represent exceptional states, not probabilities. A CLI or MCP wrapper must translate them into explicit unavailable-test statuses while retaining the original diagnostic. A downstream agent that merely compares the field with $0.05$ could otherwise describe a failed test as evidence of no effect.~\citep{exec25_rnet_biclique}

The selected randomized-tail helper counts approximate ties with $|t_r-t_0|<10^{-14}$, but separately counts all strict exceedances $t_r>t_0$. Those two sets can overlap. Its expression is
\[
 p_U=\frac{\#\{r:t_r>t_0\}+U(1+\#\{r:|t_r-t_0|<\varepsilon\})}{R+1}.
\]
Take $t_0=0$, simulated values $(10^{-15},10^{-15},1)$, $\varepsilon=10^{-14}$ and the stipulated draw $U=3/4$. There are three strict exceedances and two approximate ties, giving $p_U=21/16$. The two-sided expression $2\min(p_U,1-p_U)$ then gives $-5/8$. The nondegeneracy condition is satisfied in this authored example. This is a counterexample for the selected arithmetic contract, not a native package run. A proposed repair must partition ties and strict tails consistently and then prove or test the intended randomization law; clipping a number into $[0,1]$ is not a validity argument. Random-state identity, tolerance scaling and refusal behavior belong in the qualification suite.~\citep{exec25_rnet_biclique,exec25_rnet_choice}

\subsection{Policy targeting: numerical optimization does not restore missing graph truth}

NetworkTargeting proposes welfare-maximizing allocations under network features and treatment constraints. Its selected wrapper refuses missing treatment or primary features, but handles other missing inputs by changing the focal population and imputing values. Missing contextual covariates receive column means; missing entries of selected network matrices receive zero; affected rows become nonfocal. A zero is an asserted absence of a relationship, whereas an unknown edge is an observation problem. The resulting policy objective therefore needs its own explicit population and missingness interpretation.~\citep{exec25_rnet_target}

The inspected MILP wrapper forwards the first constraint-tolerance parameter to both tolerance arguments. The solver helper has explicit feasibility settings and a time limit when it constructs its own parameters, then returns the Gurobi result. The selected caller consumes the solution vector and objective in subsequent checks before an explicit termination-status guard in that path. An admissible policy response must distinguish a feasible incumbent, an optimal solution within tolerance, an infeasible request, a time-limited result and missing solution fields. It must also preserve focal population, budget and out-of-sample policy value. A finite assignment vector alone certifies none of those properties.~\citep{exec25_rnet_target}

Source rights also need a concrete resolution. A root GPL version-three text is present, while DESCRIPTION still contains a license placeholder; the selected dependencies include Gurobi. Those observations are not a complete distribution or entitlement decision. The proposal keeps this candidate as a policy-optimization reference pending coherent rights, solver and scientific contracts. Direct embedding or porting is deferred. The known route is to fix the target, trace the complete optimization and score paths, test explicit termination states, and compare a frozen public allocation problem against an exact small reference.~\citep{exec25_rnet_target,exec25_rnet_choice}

\subsection{Heterogeneous network effects: cluster splitting does not justify changing exposure}

NetworkCausalTree builds a discovery tree from sampled clusters and estimates effects on the remaining clusters. Its estimation path constructs the induced adjacency matrix of the retained units and recomputes neighborhood exposure there. When clusters are disconnected components, this can preserve the relevant relationships. If treatment across a discarded cross-cluster tie affects a retained unit, the induced graph changes the exposure definition. Honest sample splitting and scientifically correct exposure construction are separate requirements. The bank must supply a defensible cluster boundary or retain the appropriate background exposure.~\citep{exec25_rnet_tree}

The public argument describes assignment probability as a unit-specific vector. The inspected marginal helper nevertheless computes the chance of at least one treated neighbor using the focal unit's probability raised to its degree. For a focal probability $1/5$ and one neighbor probability $4/5$, independent assignments give
\[
 \Pr(W_i=0,G_i=1)=(1-1/5)(4/5)=16/25,
\]
while the selected expression uses $(1-1/5)(1/5)=4/25$. Equal assignment probabilities avoid this particular discrepancy. The adapter must therefore qualify heterogeneous assignment probabilities explicitly rather than infer support from a vector-shaped argument.~\citep{exec25_rnet_tree,exec25_rnet_choice}

There is a separate adjacent-pair issue to test. In a two-node connected graph with independent probability-one-half assignments, both units having own treatment one and a treated neighbor occurs exactly when both are treated, with probability $1/4$. The selected joint helper expression evaluates to $1/16$ when the shared-neighbor count is zero. The inspected variance path, however, calls its pair correction only when that count is positive: adjacent nodes with no common third neighbor are a separate dependency case. The authored helper witness and the actual pair-selection branch must both be tested; they should not be conflated into an assertion about a completed native interval. The proposal defers native admission pending exposure, marginal, joint-probability, pair-selection and tree-selection qualification.~\citep{exec25_rnet_tree,exec25_rnet_choice}

\subsection{Candidate graphs and auto-g-computation answer different uncertainty questions}

MisspecifiedInterference takes a list of candidate graphs and forms an exposure indicator that requires the same condition under every graph in that list. Its simulated probability matrices use an additive diagonal correction, and its mean routines expose both fixed-population and estimated-denominator forms. This construction makes the candidate list scientifically consequential. If one candidate leaves the focal unit isolated, an indirect-exposure condition can be impossible under that candidate. Requiring agreement with a one-edge alternative can then have probability zero. A diagonal correction of $1/(R+1)$ after zero observed occurrences cannot make that target identifiable. At $R=9$ it is numerically $1/10$, while structural support remains zero.~\citep{exec25_rnet_nmr,exec25_rnet_choice}

The route forward is to separate known design probabilities, simulation error, candidate-graph restrictions and substantive guarantees that the true exposure model is represented. An empty common event requires refusal, a changed target or justified information; it is not a numerical smoothing problem. The selected DESCRIPTION contains an unevaluated license-helper placeholder and the inspected complete repository tree has no license original. This candidate is retained for reference study, with copying or distribution deferred until rights and native behavior are resolved. No general repair for unknown bank relationships is inferred.~\citep{exec25_rnet_nmr}

autognet instead evaluates network response models through parameter draws and nested Gibbs simulations. The inspected R method selects parameter iterations by burn-in and thinning, converts adjacency identities for native calls, and constructs average, direct and spillover contrasts from distinct simulated targets. The native individual-intervention path runs a separate chain for each selected person. Its averaging mode retains sampled binary outcomes after burn-in; the nonaveraging mode returns the final conditional response probability. These are different Monte Carlo outputs even when they concern the same model.~\citep{exec25_rnet_auto}

The number of parameter draws, number of within-network simulation sweeps, initial states, retained iterations and individual intervention population all belong in the result record. More sweeps may reduce simulation error conditional on a model, but do not validate that model or certify chain mixing. The inspected source provides a bounded integration candidate after index, allocation, convergence, Monte Carlo error and dependency tests. An exact small-network enumeration can serve as the first reference before larger public comparisons.~\citep{exec25_rnet_auto,exec25_rnet_choice}

\subsection{Clustered policies: an available-group mean is not an unchanged population}

clusteredinterference supplies a group-based policy workflow with estimated nuisance parameters, treatment-vector evaluation and a sandwich variance. Its selected allocation helper samples distinct vectors with a prescribed treated count and carries a multiplicity factor into the probability integrand. Asking for all vectors sets the target number to the combinatorial support size, but still collects them through random proposals and duplicate rejection. For six equiprobable vectors the expected number of proposals to collect all six is $6\sum_{j=1}^6 1/j=147/10$. There is no deterministic upper bound on draws from that loop. A complete public test should report proposal count and incomplete work, not only the requested number of distinct vectors.~\citep{exec25_rnet_cluster,exec25_rnet_choice}

The quadrature helper integrates a single random effect over five standard deviations on either side of zero. A caught integration error or a missing integral returns numeric zero. Numerical failure and a scientifically zero assignment probability are different states and must remain distinct in the adapter. The selected target routine takes the mean of group contributions with missing values removed. For contributions $(2,\text{missing},8)$, the available-group mean is five; summing available contributions over the original three-group denominator gives $10/3$. Neither calculation supplies the unknown contribution. The distinction matters particularly because earlier inferference paths retain a different denominator convention. Group omissions must be visible before any cross-package comparison.~\citep{exec25_rnet_cluster,exec24_interference_inferference,exec25_rnet_choice}

The sandwich routine solves an assembled bread matrix and sums group-score outer products. Its inputs, derivative choices, singularity behavior and relation to the actual number of usable groups must be reviewed together. The proposal retains a conditional group-policy reference, with numerical-failure typing and population accounting as explicit admission requirements. Cluster labels cannot establish absence of cross-group interference or shared calendar shocks.~\citep{exec25_rnet_cluster}

\subsection{A public qualification portfolio with independent candidate decisions}

The first public fixtures should enumerate tiny assignment spaces and verify exposure categories, support, marginals and joint probabilities exactly. They should include asymmetric graphs, a triangle, adjacent nodes without common neighbors, heterogeneous assignment probabilities, isolated nodes and cross-cluster ties. The second group should test numerical and output contracts: empty exposures, weak instrument denominators, indefinite score kernels, near ties, failed randomization statistics, missing group contributions, quadrature errors, time-limited optimization and nested simulation retention. Every request must retain the scientific target and a typed failure state across library, CLI and MCP.~\citep{exec25_rnet_choice}

The third group concerns evidence beyond arithmetic. Match the assignment and sampling law, compare uncertainty over fresh public or generated studies, retain complete failure denominators, and measure total adapter, compute and review costs. An illustrative USD 30,000 program with 1,000 attempts, 600 completions and 400 accepted outputs costs USD 30 per attempt, USD 50 per completion and USD 75 per accepted output. These exact arithmetic comparisons are hypothetical and do not establish bank returns. A useful public result can justify later bank experiments; it cannot supply missing assignment truth or graph ownership.~\citep{exec25_rnet_choice}

The eight candidates remain separate reuse decisions. Bounded searches contain irrelevant interference applications, discovery-only repositories and uninspected tails. Other observational, simulation, spatial and experiment-design families remain open; the selected originals do not establish exhaustive R coverage. Complete delegates, native artifacts, dependency rights and scientific restrictions remain explicit. A locally unresolved license or source path does not prevent another candidate's qualification from advancing. The proposal therefore records an inconclusive coverage outcome while specifying executable next tests, without claiming complete network qualification.~\citep{exec25_rnet_choice}


% SOURCE docs/release/2026-10-29/source/docs/writing-repair/2026-10-29/src/04b976aba7e67a1e.tex CHAR 0
\section[Python network inference]{Python network inference: preserving the assignment law through learning and policy evaluation}
\label{sec:exec25-python-network}

A bank allocating interventions across connected borrowers needs more than a graph attached to an estimator. A household relationship can make an apparently untreated customer indirectly exposed; a shared employer can create common shocks; a branch allocation can induce correlated assignments; a recommendation policy can select a different population from the one used to train a model. These distinctions determine both the target and the uncertainty. The Python comparison therefore extends the existing MossSpider review with six source collections addressing different parts of the problem: ANI, GraphDML, targeted interference, CausalMotifs, CausalGraphCut and Causal-PIE. None is treated as a universal network-causal engine.~\citep{exec24_interference_mossspider,exec25_pnet_ani,exec25_pnet_dml,exec25_pnet_targeted,exec25_pnet_motifs,exec25_pnet_cut,exec25_pnet_pie}

The selected originals are pinned to commits dated 8 June 2025, 7 October 2026, 10 August 2026, 8 February 2021, 26 May 2025 and 21 March 2025 respectively. GraphDML declares version 0.1.0. Its selected commit is at 15:30 UTC on the cutoff day, or 23:30 in Asia/Shanghai, within the client cutoff. Forty-eight saved files, complete selected repository trees and inspected technical anchors support the comparison. Repository search supplies a bounded candidate list with 54 unique leads; uninspected leads and search truncation remain omissions. The calculations below are project-authored checks of specified expressions. No acquired package, upstream test, neural training routine or graph simulator was executed, and papers mentioned in source comments are not silently adopted as theorem evidence.~\citep{exec25_pnet_choice}

\subsection{A common interface must preserve distinct network targets}

Write $D_i$ for the treatment assigned to node $i$, $A$ for the declared network, and $e_i(D,A)$ for its exposure. The response may be represented as $Y_i(d,e)$ only under a substantive restriction connecting the full assignment vector to the chosen exposure. A graph feature learned for prediction does not establish that restriction. A design-based contrast requires assignment probabilities for exposure events; a partially linear model instead imposes a response structure; a policy evaluator needs the law of the proposed assignments and the relevant time horizon. These objects should not be collapsed into a generic effect field.

For banking use, the request should therefore bind the observation time, node identifiers, arrow direction, edge weights, treatment assignment law, exposure operator, analysis population, target policy and uncertainty law. A sum of neighbors' treatments measures a different intervention from their mean, especially when degree varies. An isolated borrower needs an explicit exposure convention. A missing relationship is not automatically an absent relationship. The proposed CLI and MCP result must carry these decisions and any refusal, so downstream agents cannot interpret an output number without its scientific contract. The package comparisons provide concrete tests of this proposed interface; they do not establish that the interface is implemented.~\citep{exec25_pnet_choice}

\subsection{ANI: uncertainty must not depend on the order of requested contrasts}

ANI constructs an inverse-probability exposure contrast from supplied indicators and probabilities. Its population mask selects the contributors, while excluded score entries receive a sentinel that is not intended for estimation. The standard-error routine computes distances in the full undirected, unweighted graph and assumes nodes are numbered consistently with score positions. Thus subsetting outcomes does not silently mean recomputing distances on an induced graph. The explicit zero-bandwidth branch uses the population-standard-deviation convention divided by the square root of the selected count.~\citep{exec25_pnet_ani}

For a positive bandwidth, let $z$ be the centered score vector over the selected population of size $n$ and let $K_{ij}=1\{d_A(i,j)\leq b\}$. The selected variance of the mean is $z'Kz/n^2$. When the intermediate quadratic is nonpositive, the code constructs a replacement from overlapping neighborhoods. If $B_{ij}=1\{d_A(i,j)\leq b\}$ and $r_i=\sum_jB_{ij}$, that replacement is $P=\widetilde B\widetilde B'$ with $\widetilde B_{ij}=B_{ij}/\sqrt{r_i}$. Its positive-semidefinite form is apparent from the product, but the implementation places the replacement into a weights variable shared across the loop over contrasts. Subsequent contrasts use that replacement even when their original quadratic would have been positive. The fallback flag records whether the current contrast triggered the branch, not necessarily which kernel it inherited.~\citep{exec25_pnet_ani}

An authored three-node path illustrates the interface risk. Set $b=1$ and
\[
 K=\begin{pmatrix}1&1&0\\1&1&1\\0&1&1\end{pmatrix},
 \qquad u=(1,-2,1)',\qquad v=(1,0,-1)'.
\]
Both scores are centered. Their original variance quadratics divided by nine are $-2/9$ and $2/9$. Processing $u$ first activates the overlap kernel. The resulting variance for $u$ is $(7-16/\sqrt6)/9$, and $v$ then receives $1/9$. Processing $v$ first instead gives it $2/9$. The latter contrast can therefore change solely because another requested contrast preceded it. This is a source-expression witness, not a reported run of ANI on a fitted banking sample.~\citep{exec25_pnet_choice}

ANI is a conditional reuse candidate for design-based network uncertainty. The adapter should preserve the original kernel, allocate any fallback locally to each contrast, report the kernel actually used, and test permutation invariance of a batch of contrasts. It must also test disconnected and edgeless graphs, singleton components, node ordering, structural support, nonfinite scores and subsampling. These engineering tests do not prove the sampling assumptions needed for a network central-limit approximation. The authentic BSD three-clause license is a useful starting point; installed dependencies and distribution obligations remain a separate review.~\citep{exec25_pnet_ani,exec25_pnet_choice}

\subsection{GraphDML: useful graph-aware construction with population and covariance gates}

GraphDML makes several interface distinctions explicit. The data container checks covariate finiteness and graph dimensions, warns when dropping self-loops, rejects a nonsymmetric graph unless direction is declared, and checks node identifiers. Its exposure operators distinguish unweighted sums, unweighted means, weighted sums and means, and shortest-path distance exactly two. The two-hop implementation removes self and direct-neighbor entries from the thresholded matrix product. The mean operator leaves isolated rows at zero. Those conventions are valuable adapter contracts, provided the bank accepts their causal meaning.~\citep{exec25_pnet_dml}

The estimation path forms dependency sets from own and exposure-row support. For a held-out focal set, buffered nuisance training excludes treatment observations in the union of its dependency sets and excludes outcome observations whose own dependency sets overlap that union. The default final covariance is an empirical sandwich without a supplied group vector. This is not evidence that the package lacks a cluster option: a separate cluster-summary method aggregates scores by a supplied group and explicitly requires one repetition. Disjoint graph supports alone do not verify conditional independence of primitive disturbances. A common employer shock can violate that interpretation despite distant graph nodes; the scientific model and a justified cluster law remain necessary.~\citep{exec25_pnet_dml}

The population contract also needs attention. When the focal set is a named selection rule, the declared estimation nodes are supplied as candidates. When explicit focal indices are supplied, the selected branch checks index range and disjoint dependency supports but returns the indices without intersecting them with the declared estimation population. Total-effect weights subsequently use that declared population. An adapter should reject incompatible index sets or explicitly justify transport of the fitted coefficients; it should not silently describe coefficients from one population as estimated in another. This is a branch-level finding, not a claim that every fit takes the incompatible path.~\citep{exec25_pnet_dml}

The total effect has a precise model interpretation. For the proposed partially linear response
\[
 Y_i=g(X_i)+\theta D_i+\sum_{k=1}^q\beta_k(E_kD)_i+\varepsilon_i,
\]
with appropriately identified common coefficients, changing every treatment from zero to one gives $\theta+\sum_k\beta_k\overline{E_k\mathbf1}$ over the declared population. This follows by subtraction under the stated model; it does not identify an unrestricted heterogeneous policy effect. The source computes these weights and uses their covariance quadratic for total-effect uncertainty.~\citep{exec25_pnet_dml,exec25_pnet_choice}

For repeated fits, the selected helper takes a coordinatewise median of coefficients and then an elementwise median of covariance matrices adjusted by coefficient-displacement outer products. Elementwise medians need not preserve positive semidefiniteness. Consider identical coefficient vectors and three positive-definite covariance inputs
\[
 C_1=\begin{pmatrix}1&-2\\-2&5\end{pmatrix},\quad
 C_2=\begin{pmatrix}5&-2\\-2&1\end{pmatrix},\quad
 C_3=\begin{pmatrix}1&0\\0&1\end{pmatrix}.
\]
Each has positive leading entry and determinant one. Their elementwise median is $M=\left(\begin{smallmatrix}1&-2\\-2&1\end{smallmatrix}\right)$, with determinant $-3$. For total-effect weights $(1,1)$, the variance quadratic is $-2$. This exact helper-input counterexample does not establish that a particular fitted repetition produces these matrices. It establishes that entrywise aggregation alone cannot certify a covariance suitable for every linear contrast. Merely clipping a negative variance does not establish correct coverage.~\citep{exec25_pnet_choice}

The proposed decision is conditional adoption through a narrow adapter: retain graph-aware splits, exposure operators, diagnostics and explicit cluster summaries, while requiring population compatibility and contrast-valid covariance aggregation. Public trials should include common shocks, near-singular moments, repeated splitting, isolated nodes, directed graphs, exact two-hop semantics and changed node ordering. The MIT original and dependency declarations permit a concrete rights review; they are not complete artifact clearance. No native capability is admitted from these static findings.~\citep{exec25_pnet_dml,exec25_pnet_choice}

\subsection{Targeted interference and CausalMotifs: learning the representation is not the entire inference problem}

The targeted-interference implementation combines a graph encoder, separate outcome heads, a treatment-probability head, a neighborhood-density head and spline fluctuations. The selected neighborhood mean is $AD/(A\mathbf1)$ with no isolate guard in that expression. The treatment head is trained with binary cross-entropy, while the fluctuation correction uses the product of its output and the neighborhood-density output plus a small numerical stabilizer. The treatment of control targets and the meaning of that factorization need a method-specific derivation and known-law comparison. The inspected expression alone does not justify declaring the entire method incorrect.~\citep{exec25_pnet_targeted}

The selected code also includes CUDA-specific allocations and evaluates supplied test potential-outcome truths during training methods. The inspected top-level loop uses configured epoch counts; it does not establish that the optimizer selects a checkpoint using test truth. Nevertheless, a proposed public experiment should keep final evaluation oracles inaccessible until the model and hyperparameters are frozen. Synthetic potential outcomes can supply that oracle, whereas observed bank data generally cannot. The conditional recommendation is to study the architecture as a research reference and defer direct adoption until the target derivation, isolate handling, device portability, held-out evaluation and distribution rights are resolved. No license-named original appears in the selected complete tree.~\citep{exec25_pnet_targeted,exec25_pnet_choice}

CausalMotifs learns partitions using network-motif features and simulated exposure probabilities. Its constructor randomly divides rows into training and estimation subsets. That split is useful for separating partition selection from estimation on rows, but does not itself make their network disturbances independent. The nonseparate validation branch removes rows with zero empirical probabilities and fits weighted least squares using the default fit covariance. The separate H\'ajek helper contains individual variance and covariance terms rather than explicit pairwise exposure covariances. Appropriate uncertainty therefore needs a declared assignment/dependence law and a target population that records the filtering.~\citep{exec25_pnet_motifs}

Two selected support checks make small, deterministic tests especially valuable. The branch without an own-treatment restriction reads the stored probability collection rather than the passed collection. The own-treatment branch adds a list of rule arrays to the treatment array before comparing with the target value; it does not append a parenthesized treatment-condition array to the rule list. In the scalar arithmetic analogue, a true rule and treatment one produce $(1+1)==1$, which is false although their conjunction is true. Separately, the caller computes left-control and right-control support checks but repeats the right-control term in its rejection predicate. Array shapes and the exact caller path must be qualified in native tests before interpreting the practical effect.~\citep{exec25_pnet_motifs,exec25_pnet_choice}

The example generator requires equal care. It uses a very large synthetic graph, clustered labels, and an index-distance shortcut to judge selected neighbor-neighbor relationships. Its subsequent probability simulation rotates cluster alignment, unlike the initial assignment expression. A bank should not port those illustrative shortcuts as general graph or assignment semantics. The MIT source is a conditional research/adapter candidate after explicit Boolean-rule tests, assignment-law reconciliation, support accounting and dependent uncertainty. Native qualification remains open.~\citep{exec25_pnet_motifs,exec25_pnet_choice}

\subsection{CausalGraphCut: inspect the custom delegate before judging the interface}

CausalGraphCut addresses experimental allocation rather than simply estimating an effect after assignments have been made. It constructs a covariance-weighted graph objective and delegates spectral partitioning to a local class. That class directly eigendecomposes its input for the selected discretization path. It would therefore be wrong to infer a standard-library affinity-interface error merely from the public class name: the actual delegate has different semantics. The proposed reuse is an allocation-design research reference whose objective must match the bank's assignment probability, network and outcome dependence.~\citep{exec25_pnet_cut}

The inspected source exposes separate implementation gates. The default maximum cluster count is a rounded power of the node count cast to a signed eight-bit integer. For 4,096 nodes the unwrapped value is 256, whose eight-bit wrapping representation is zero. This is an authored arithmetic witness; installed NumPy conversion behavior remains a native test. The online routine resets batch size to five, uses complete batches only, and depends on a graph variable populated by the main example rather than an explicit graph argument. With seven requested samples, its batch arithmetic schedules five and omits two. With only one complete batch, the returned cluster variable is not assigned in the initial branch. Larger studies also need to retain the full allocation and adaptation history.~\citep{exec25_pnet_cut,exec25_pnet_choice}

The objective list is sorted in descending order and the online caller selects its last element, which correctly corresponds to the smallest listed objective; this is not a sorting defect. Other scope distinctions remain: the selected objective helpers use half-probability constants, while design classes accept a probability parameter and use different neighborhood or cluster probability exponents. Generalizing the objective to arbitrary assignment probabilities requires a derivation, not just changing a parameter. Direct adoption is deferred pending explicit graph arguments, safe integer handling, sample accounting, typed return paths, adaptive-design validation and rights. No license-named original appears in the selected complete tree.~\citep{exec25_pnet_cut,exec25_pnet_choice}

\subsection{Causal-PIE: an embedding match is not the assignment probability}

Causal-PIE connects learned neighborhood representations to policy evaluation. In its one-stage path, records are selected when own action matches the target and a learned embedding is sufficiently close. The relative criterion uses a strict inequality between root-mean-square distance and a threshold times the root-mean-square reference embedding. Identical zero embeddings fail this test because $0<0$ is false. Moreover, filtering and probability calculation reverse which embedding supplies the reference scale. An adapter must specify a symmetric or intentionally asymmetric equivalence rule, zero-scale behavior and the scientific adequacy of the learned exposure representation.~\citep{exec25_pnet_pie}

The probability calculation offers Monte Carlo and traversal modes. Monte Carlo draws independent neighbor actions using the supplied Bernoulli probability $p$. Traversal enumerates binary configurations but its active expression weights them uniformly. After clipping the neighborhood result, the code multiplies by $p$ regardless of whether the target own action is zero or one. These choices coincide with some balanced designs, but not arbitrary assignment laws.~\citep{exec25_pnet_pie}

For a transparent authored example, suppose one neighbor is treated independently with $p=1/5$, the target own action is zero, and the target neighbor action is one. Assume a nonzero injective embedding, no added noise and an inactive clipping bound. The correct joint assignment probability is $(1-p)p=4/25$. Uniform traversal followed by multiplication by $p$ gives $1/10$; the Monte Carlo neighborhood limit followed by that same own-action factor gives $1/25$. The experiment should test both probability modes at balanced and unbalanced assignments and for both own actions. A floor on the estimated probability cannot repair the wrong assignment law.~\citep{exec25_pnet_choice}

A second composed-path witness concerns empty support. When no record passes the filter, the selected routine returns record index zero with probability one. The caller then computes a weighted residual correction using the returned record and divides by the full number of node-time observations. If that first residual were six and the full count ten, the correction would be $3/5$ rather than the zero sum over an empty selection. This is hypothetical arithmetic, not a measured dataset. The appropriate product response should report empty selected support, and either return the explicitly defined plugin-only quantity with its limitations or refuse the requested correction. A numerical sentinel must not become an observed policy match.~\citep{exec25_pnet_pie,exec25_pnet_choice}

The multi-stage implementation poses different questions. It combines a fitted value function, an RKHS-trained density ratio, discounted next values and an initial-value plugin normalized by $1-\gamma$. Its selected constructor assigns the same trajectory indices to training and test collections. The residual uses the node, trajectory and $T-1$ transition count, so finite-horizon terminal handling and stationary-policy interpretation need separate qualification. The RKHS bandwidth is a median embedding distance, creating a zero-bandwidth gate when embeddings collapse. These are concrete proposed tests; neither the source's double-robust naming nor an accurate fitted reward model establishes valid policy uncertainty.~\citep{exec25_pnet_pie}

The candidate is retained as a policy-evaluation research reference, with direct adoption deferred until probability-law weighting, empty support, matching scale, trajectory separation, horizon semantics, collapsed representations and rights are resolved. The selected complete tree contains no license-named original. Public synthetic trajectories with known intervention values can test these contracts before any expensive bank study. Unknown confounding and an inadequate exposure map still require scientific restrictions or more information.~\citep{exec25_pnet_choice}

\subsection{Integration decisions and a bounded public qualification sequence}

The six decisions are independent. ANI and GraphDML supply concrete adapter candidates after uncertainty and population repairs; CausalMotifs supplies a conditional partition-learning reference after support and dependence checks. The neural targeted, allocation-design and policy-evaluation repositories supply useful architecture and experiment references, with direct adoption deferred until their stated rights and technical gates are resolved. A port to another language should preserve the declared mathematical contract, not reproduce an accidental sentinel or change the population silently. Direct reuse is preferred only when the actual artifact, license and scientific regime meet that contract.~\citep{exec25_pnet_choice}

The first public qualification stage should enumerate tiny assignment spaces and test exposure orientation, isolate behavior, exact joint support, own-action probabilities and policy denominators. The second should examine covariance aggregation, batch-order invariance, common shocks and graph-aware splitting on frozen known-law data. The third should freeze model selection and evaluate all requests, including failures, empty support and nonfinite outputs, on held-out graph and trajectory families. A comparison records identification assumptions, sampling and Monte Carlo uncertainty, runtime, complete costs, refusal accuracy and target-population coverage separately. It cannot manufacture missing bank assignment truth or prove universal validity under arbitrary hidden network dependence.~\citep{exec25_pnet_choice}

For scale, a hypothetical USD 30,000 study with 1,000 attempted requests, 600 completed requests and 400 scientifically accepted outputs costs USD 30 per attempt, USD 50 per completion or USD 75 per accepted output. These are three distinct denominators, not observed return on investment. Bounded search tails, complete dependency/artifact review, native reproduction and independent acceptance remain open. The selected primary sources and exact authored witnesses establish a reviewable route to public qualification, not complete Python network qualification.~\citep{exec25_pnet_choice}


% SOURCE docs/release/2026-10-29/source/docs/writing-repair/2026-10-29/src/a6a9c3beac1b082c.tex CHAR 0
\section[JVM network contracts]{JVM network and simulation components: making exposure and intervention clocks explicit}
\label{sec:exec25-java-network}

A bank constructing borrower-network experiments on the JVM must separate the graph representation from the causal experiment. The graph can encode households, suppliers, branches or common employers; its edges may mean influence, contractual exposure or merely similarity. A simulator can generate outcomes after an intervention has been specified, but its output does not identify that intervention from observed bank data. This comparison evaluates JGraphT, GraphStream, MASON and Repast Simphony as components for graph construction and public experiments. Admission as a component is separate from qualification of a causal estimator.~\citep{exec25_jnet_jgrapht,exec25_jnet_graphstream,exec25_jnet_mason,exec25_jnet_repast}

The inspected revisions precede the 7 October 2026 client cutoff. JGraphT's selected tree declares 1.6.0-SNAPSHOT and Java 21; the release snippet in its README is not the selected build. GraphStream declares 2.1.0-SNAPSHOT in the selected November 2021 tree. The MASON review binds the July 2026 source revision, while the Repast core manifest identifies 2.12.0 at a September 2026 revision. Fifty technical, license and configuration originals plus two candidate-screening originals were saved and authenticated. Eight bounded repository queries return 71 distinct leads. Some query results have uninspected tails, and the returned hits include computing, radio, genetics and distributed-system meanings of interference or causality. These results cannot establish absence of Java network-causal software.~\citep{exec25_jnet_choice}

One seemingly relevant partial-interference repository describes colocated server workloads, resource metrics and an SLA scheduler in its original README. It is not selected as evidence of a borrower-spillover estimator. A solar-innovation search lead has an empty root README and remains unresolved without technical inspection. The distinction prevents a search label from becoming a capability claim. No downloaded code, native JVM artifact, upstream test or simulator was executed; no theorem or benchmark is imported from papers mentioned in source comments.~\citep{exec25_jnet_choice}

\subsection{JGraphT: neighbor lists, graph distances and live subgraphs}

JGraphT provides distinct operations that a causal adapter should preserve. Its neighbor list traverses all incident edges and can include a vertex more than once in a multigraph; its neighbor set removes duplicate vertices. Predecessor and successor lists instead traverse incoming and outgoing edges. These are useful alternatives, not interchangeable encodings of a bank exposure. Whether repeated transactions count as repeated influence, summed exposure weight or a single relationship is a scientific decision.~\citep{exec25_jnet_jgrapht}

Consider the authored directed multigraph with two edges $a\to b$ and one edge $b\to c$, with treatment vector $(D_a,D_b,D_c)=(1,0,0)$. For node $b$, the mean treatment over the incoming list is one. Over the full incident-neighbor list it is $2/3$, because $a$ appears twice and $c$ once. Over the unique adjacent-vertex set it is $1/2$. A downstream model that sees only a field named neighbor exposure cannot distinguish these quantities. The proposed adapter should record direction, multiplicity, weighting, self-edge handling and normalization with the exposure operator.~\citep{exec25_jnet_choice}

The breadth-first shortest-path implementation traverses outgoing edges and increments distance by one. It therefore computes hops, regardless of stored edge weight. Its returned single-source object retains those hop distances; when reconstructing a path, the delegate sums the actual weights of selected edges. In an authored triangle with weights $w_{ab}=10$, $w_{ac}=w_{cb}=1$, the direct path has one hop and reconstructed weight ten, while the minimum-weight route has two edges and total weight two. The hop-distance query, the weight of the hop-selected path and the minimum-weight distance are three different quantities. An unreachable target has an infinite stored distance and no reconstructed path; it must not be silently converted to a zero-distance neighbor.~\citep{exec25_jnet_jgrapht,exec25_jnet_choice}

Subsetting also changes the question. In the path $a\to b\to c$, the full-graph distance from $a$ to $c$ is two. The induced graph on $\{a,c\}$ has no such path. Filtering the analysis population while retaining background network paths is therefore different from filtering the graph itself. This distinction matters when a bank observes outcomes only for a subset of borrowers but influence passes through other entities.~\citep{exec25_jnet_choice}

The selected subgraph class is a view over a backing graph. Its documentation and implementation distinguish listener-aware topology updates from a non-listenable base. Its edge-weight setter explicitly delegates to the backing graph. General wording that subgraph modifications do not affect the base must consequently be read as a topology statement, not a guarantee that all attributes are isolated. A scenario that changes edge weights in a subgraph can change the backing graph used by another scenario. The proposed contract uses immutable input identifiers and a deliberate deep-copy or versioned-view policy, with tests of both topology and attributes.~\citep{exec25_jnet_jgrapht}

JGraphT is a conditional direct-integration candidate for graph operators. The source supports explicit seeds or a supplied generator for its selected preferential-attachment generator, but repeatable generation still depends on vertex construction, iteration order, algorithm version and the graph's admissible edges. Synthetic topology is not authentic bank structure. The dual LGPL 2.1-or-later or EPL 2.0 declaration and complete local license texts support a concrete distribution review, while the selected snapshot, Java requirement and transitive components still need artifact-specific qualification. Porting should be considered only where the target runtime or numerical contract justifies the additional maintenance burden.~\citep{exec25_jnet_jgrapht,exec25_jnet_choice}

\subsection{GraphStream: ingestion policy changes the observed population}

GraphStream provides a different component: a graph that evolves through events. Its default edge-addition overload creates an undirected edge, while an explicit directed addition uses the specified source-to-target orientation. The DGS parser maps forward and reverse arrow tokens to corresponding directed events, and a step token becomes a numerical graph-step event. That step is not, by itself, the bank's event time, ingestion time or availability time. Those clocks require explicit fields and a documented conversion.~\citep{exec25_jnet_graphstream}

Error behavior is controlled by strict checking and automatic creation. A duplicate identifier may throw, return an existing compatible edge, or return null for incompatible endpoints. Missing endpoints may throw, produce null, or be created automatically. Edge rejection by a node similarly can throw or return null, with rollback of the first endpoint when necessary. These paths affect the observation contract: automatic creation can enlarge the population, while a discarded null can remove a relationship without telling the analyst. A causal ingestion adapter should return a receipt for each requested node or edge, including accepted, deduplicated, created, rejected and unresolved outcomes.~\citep{exec25_jnet_graphstream}

SingleGraph and MultiGraph differ at their node delegates. The selected SingleNode stores incoming and outgoing slots per opposite node and rejects conflicting additions; MultiNode retains a list of parallel edges. A read-only collection wrapper over that list does not freeze the underlying edge objects. A bank replay should therefore bind both the ordered event sequence and the resulting immutable graph snapshot, rather than retain a mutable collection reference as historical evidence. It should test repeated identifiers, reverse arrows, parallel obligations, unknown nodes and late-arriving attributes.~\citep{exec25_jnet_graphstream}

The proposed decision is conditional integration for dynamic ingestion and visualization, with causal exposure construction in a separately specified layer. The source's dual CeCILL-C or LGPL 3-or-later declaration requires a distribution choice and dependency review; it is not an unrestricted artifact clearance. No benefit follows from porting event handling unless the destination must preserve the same failure, ordering and time semantics. An agent choosing a permissive ingestion mode should not silently change the target population.~\citep{exec25_jnet_graphstream,exec25_jnet_choice}

\subsection{MASON: simulation state, update order and checkpoint identity}

MASON's network stores arbitrary node objects and edge information. Incoming and outgoing edge bags are internal mutable structures, and adjacency structures must be regenerated after topology changes. Its simple adjacency matrix is indexed from source to target and stores a single selected edge for a pair; a separate multigraph method is available when all parallel edges are needed. The selected weight getter returns a numeric or Valuable information object's value, and otherwise returns one. Its weight setter replaces the information object with a numeric value. A bank adapter should reject an unexplained unit-weight fallback where an edge actually represents a missing monetary amount or an uncertain relationship.~\citep{exec25_jnet_mason}

Copying the network does not automatically isolate a counterfactual scenario. The selected copy routine reuses node objects and passes the same edge-information references into new edge wrappers. A mutable customer state or edge-information object can therefore be shared across copies. The proposed scenario runner must define which state is copied, which is immutable, and which exogenous draws are intentionally coupled. It should test that changing a treatment or attribute in one scenario does not unintentionally mutate another. A full-state serialization checkpoint and a graph-wrapper copy solve different problems.~\citep{exec25_jnet_mason}

The scheduler gives equally consequential semantics. Lower integer orderings run first; equal-time, equal-order substeps are shuffled and then invoked sequentially. A newly scheduled event at the current time advances to the next representable time step. Thus a group of events with the same nominal timestamp is not necessarily a simultaneous update from the previous state. A simulator needs an explicit read-old/write-new protocol when that is the intended causal mechanism.~\citep{exec25_jnet_mason}

An authored two-node example isolates this point. Start from $(x_1,x_2)=(1,0)$ and define each update as copying the other node's current value. A simultaneous update from the old state gives $(0,1)$, with mean $1/2$. Sequential order $(1,2)$ gives $(0,0)$; order $(2,1)$ gives $(1,1)$. If the two sequential orders are equally likely, the mean is still $1/2$, but the variance of the resulting network mean is $1/4$, compared with zero for the deterministic simultaneous update. Agreement in an average therefore does not establish agreement in the transition law or its uncertainty. These are calculations for authored updates, not native MASON output.~\citep{exec25_jnet_choice}

Random-state identity must also be explicit. The selected seed setter narrows the supplied long integer to 32 bits, so seeds seven and $7+2^{32}$ become the same effective seed. Starting a simulation primes the generator and resets the schedule. Restoring a serialized state instead invokes a wake-from-checkpoint hook. The stream checkpoint API throws on failure, whereas convenience file wrappers print exceptions and return null. A runner should distinguish restart-from-seed from continuation-from-state and type checkpoint failure rather than interpret null as a successful empty result.~\citep{exec25_jnet_mason,exec25_jnet_choice}

The parallel-sequence documentation assigns synchronization of model data and the shared random generator to the caller. The selected RandomSequence has a switch for synchronized draws, and the fast random generator itself is unsynchronized. Locking can prevent a data race without fixing which thread consumes which random draw. Reproducible coupled experiments therefore need task-specific random streams or a deterministic allocation of draws, together with a declared update schedule and thread model. No native replay or throughput claim is made here.~\citep{exec25_jnet_mason}

MASON is a conditional simulation reference for public known-law studies. Its principal license is AFL 3.0 with named component exceptions, including the random generator's BSD terms. A proposed integration should isolate the headless simulation subset, inspect its actual distribution and test copied state, seeded starts, checkpoints, interrupted runs and independent replicates. A simulated policy result is conditional on the authored mechanism, not validation of a bank mechanism.~\citep{exec25_jnet_mason,exec25_jnet_choice}

\subsection{Repast Simphony: network direction, priority conventions and named streams}

Repast's network builder chooses directed or undirected implementations backed by ordered sparse multigraphs. Its generic adjacency method takes the union of predecessors and successors in a set, removing duplicate adjacent vertices. Separate in-edge and out-edge accessors preserve the directional distinction. Edge objects carry a mutable weight, defaulting to one. The same exposure questions therefore recur: a generic adjacent-node average cannot silently replace an incoming, weighted or transaction-multiplicity exposure.~\citep{exec25_jnet_repast}

The scheduler uses a priority convention that differs from MASON's integer ordering: larger numeric priorities execute earlier. It has explicit first, last and random-priority tokens. The inspected group delegate first establishes an order-index sequence and then intersperses random-priority actions using the random registry, or shuffles them when there are no numeric-priority actions. A port that copies numeric order values from one framework into another can reverse the intended chronology. The proposed conversion should map semantic phases, such as intervention assignment, response update and observation, before mapping numeric priorities.~\citep{exec25_jnet_repast}

The default random registry begins with a timestamp-derived integer seed. Named generators can be registered with explicit seeds. The inspected set-seed implementation invalidates default distributions, while the broader reset clears named distributions and generators and recreates the default from the clock. These distinctions are more specific than a generic statement that resetting randomness makes an experiment reproducible. A recorded seed alone also cannot establish independence of replicates or equivalence of conditional assignment laws. The runner must bind the effective generator state, name, version, draw allocation and reset path.~\citep{exec25_jnet_repast}

Repast is a conditional alternative for a richer agent-based experiment environment. The core source has BSD-style redistribution conditions and names separately licensed dependencies. The selected core manifest identifies version 2.12.0, but that does not qualify every plugin or installed runtime. The proposed choice between Repast and MASON should compare required scheduling, context, data-output and deployment features on the same public experiment, including failed requests and restart costs. A clean API response is not an identified network effect.~\citep{exec25_jnet_repast,exec25_jnet_choice}

\subsection{From simulation components to a defensible causal experiment}

To connect these components to the proposal, let $A$ denote a frozen graph, $U$ exogenous draws, $D$ an assignment and $F$ a specified state-transition system. A public simulator evaluates the consequences of $(A,U,D,F)$ under declared scheduling and observation clocks. It does not discover $F$ or establish that the observed bank assignment is independent of latent risk. Such public experiments can nevertheless test software semantics and uncertainty in regimes where the truth is known. That is the appropriate inexpensive precursor to an authentic bank study.~\citep{exec25_jnet_choice}

For two connected borrowers, independent Bernoulli assignment with probability $p$ supports all four assignments with probabilities $(1-p)^2$, $p(1-p)$, $(1-p)p$ and $p^2$. Assigning them together as one cluster supports only $(0,0)$ and $(1,1)$, with probabilities $1-p$ and $p$. At $p=1/5$, the mixed assignments have probability $4/25$ each under the independent design and zero under the single-cluster design. Neither a shared seed nor a common graph object can make those laws interchangeable. Exposure-probability and refusal tests should therefore begin with exact enumeration.~\citep{exec25_jnet_choice}

The integration plan has four independent component decisions. JGraphT supplies conditional graph operators; GraphStream supplies conditional event ingestion; MASON and Repast supply alternative public experiment environments. Direct use should preserve source contracts where they meet the scientific specification. A language port is justified only after the graph, clock, random-stream and failure contracts have an executable reference and the target runtime offers a concrete benefit. Components unsuitable for a particular target should be rejected for that use with the target-specific reason, rather than described as universally deficient.~\citep{exec25_jnet_choice}

The proposed CLI and MCP request should bind source revision and artifact, graph identity and node ordering, edge orientation and multiplicity, event and availability clocks, target population, assignment law, exposure function, update phases, effective random streams and checkpoint provenance. Its result should retain scientific assumptions, typed failures and complete costs. Small public tests can then compare directed versus incident exposure, weighted versus hop paths, induced versus background graphs, live-view versus copied state, synchronous versus sequential updates, seed collisions, restart semantics and assignment support. Native tests, full dependency rights and causal estimator qualification remain separate obligations.~\citep{exec25_jnet_choice}

Using the same hypothetical USD 30,000 study budget, 1,000 attempted requests, 600 completions and 400 accepted outputs imply costs of USD 30 per attempt, USD 50 per completion and USD 75 per accepted output. These denominators expose the economic consequence of failure and refusal without claiming measured returns. The source review identifies concrete routes to public qualification while retaining search tails, unknown bank structure and unrestricted dependence as unresolved. It does not establish complete JVM network qualification.~\citep{exec25_jnet_choice}


% SOURCE docs/release/2026-10-29/source/docs/writing-repair/2026-10-29/src/1979e0a1a216a89e.tex CHAR 0
\section[C++ network experiments]{C++ network experiments: assignment probabilities, graph identity and distributed state}
\label{sec:exec25-cxx-network}

A bank evaluating interventions across connected borrowers needs more than a fast graph engine. Its experiment must distinguish the assignment law, the exposure event that makes an observed outcome informative about a policy, and the numerical representation used to compute that event. A program may preserve adjacency while changing weights, reuse a random stream across workers, or return a number when one exposure arm is empty. Each change can alter the scientific result. This comparison examines the RGCR research implementation alongside NetworKit, SNAP, Repast HPC and the C-native igraph library as an interoperability alternative. The first supplies a specific experimental-design implementation; the others supply components whose scientific role must be declared.~\citep{exec25_cnet_rgcr,exec25_cnet_networkit,exec25_cnet_snap,exec25_cnet_repast,exec25_cnet_igraph}

The comparison binds 64 selected software, license and configuration originals at revisions before the 7 October 2026 cutoff. Six bounded repository queries produce 52 distinct leads, with unreturned search tails and uninspected candidates. They cannot establish complete C++ coverage or absence of other causal software. NetworKit's selected metadata says 11.2.1, but its source build sets C++20 and CMake 3.14; those exact requirements take precedence over a less demanding README snippet for this revision. RGCR's build instead references C++11 and a user-specific SNAP 4.1 installation. The independently inspected SNAP revision is not evidence that this older dependency has been reproduced.~\citep{exec25_cnet_choice,exec25_cnet_networkit,exec25_cnet_rgcr}

The methodological discussion uses the genuine 59-page arXiv version of Ugander and Yin, submitted on 4 September 2020, whose title page bears 7 September 2020. The later publisher edition was not inspected. Version identity matters: the argument here cites the selected author text and distinguishes its statements from the project's finite calculations and the implementation's conventions. No acquired program, simulator, native library or upstream test was executed.~\citep{exec25_cnet_paper,exec25_cnet_choice}

\subsection{Randomized clustering changes both design and analysis}

Let $G=(V,E)$ be a fixed undirected, unweighted interference graph with $n=|V|$ borrowers. For binary assignment vector $Z$, let $Y_i(Z)$ be borrower $i$'s potential outcome and define the global average treatment effect
\[
 \tau=\frac{1}{n}\sum_{i=1}^n\{Y_i(\mathbf{1})-Y_i(\mathbf{0})\}.
\]
The full-neighborhood exposure event $E_i^a$ occurs when borrower $i$ and every neighbor receive assignment $a\in\{0,1\}$. Using the observed outcome on that event requires the substantive assumption that it equals $Y_i(\mathbf{a})$. An adjacency file cannot establish this exposure model. A graph of transactions, for example, need not include all shared employers, household obligations or unobserved channels relevant to a servicing response.~\citep{exec25_cnet_paper,exec25_cnet_choice}

In randomized graph cluster randomization, a partition $C$ is sampled from a declared distribution $P$, followed by treatment assignment at the cluster level. The analysis uses the marginal exposure probability
\[
 \pi_i^a=\Pr(E_i^a\mid P)=\sum_c P(c)\Pr(E_i^a\mid C=c),
 \qquad
 \widehat\tau_{\mathrm{HT}}=
 \frac{1}{n}\sum_i\left\{
 \frac{\mathbf{1}\{E_i^1\}Y_i(\mathbf{1})}{\pi_i^1}
 -\frac{\mathbf{1}\{E_i^0\}Y_i(\mathbf{0})}{\pi_i^0}
 \right\}.
\]
The selected paper derives this marginalization explicitly. Conditioning only on the realized partition gives a different weighting scheme. Likewise, averaging inverse conditional probabilities does not compute the inverse marginal probability. If two equally likely partitions give exposure probabilities $1/4$ and $1/2$, the marginal probability is $3/8$ and its inverse is $8/3$, whereas the average of the two inverse probabilities is three. The adapter should bind the partition generator and its distribution alongside the realized partition.~\citep{exec25_cnet_paper,exec25_cnet_choice}

The source implements several partition rules, including neighborhood maxima and randomized nets. For independent cluster assignment it constructs each borrower's set of adjacent clusters, including its own, and evaluates exposure probability as $2^{-k_i}$, where $k_i$ is that set's size. It also accumulates same-arm and opposite-arm joint exposure probabilities. This is substantive support for interference-aware design analysis. It is specifically tied to the inspected half-probability assignment: the treatment simulator uses a Bernoulli distribution with parameter $1/2$. Changing only a user-facing field to $p=1/5$ would not change these expressions. For a neighborhood touching two independently assigned clusters, treatment and control exposure probabilities would be $1/25$ and $16/25$, respectively, rather than $1/4$ for each.~\citep{exec25_cnet_rgcr,exec25_cnet_choice}

The implementation's complete-randomization mode is also specific. It orders clusters by decreasing size, pairs them, and gives each pair opposite assignments; an unpaired final cluster receives its own assignment. If a borrower's neighborhood intersects both members of a pair, full exposure is impossible under that partition. The helper returns zero for this conflict, and the analysis records excluded target contributions. This does not restore positivity. Equal-size clusters have no explicit secondary key in the selected comparator, so the realized pairing and its provenance should be saved instead of reconstructed from sizes alone. The declared law is not arbitrary complete randomization or a general bank budget-constrained design.~\citep{exec25_cnet_rgcr,exec25_cnet_choice}

The paper offers a route to controlling rare exposure under its selected random partitions. For its randomized 3-net or 1-hop-max constructions, the reported treatment-exposure lower bound is $p/|B_2(i)|$, where $B_2(i)$ is the set of vertices within two hops, including $i$. Its Monte Carlo probability estimator averages conditional probabilities over sampled partitions; the stated relative mean-square error bound is $|B_2(i)|/(Kp)$ for $K$ samples. The paper also describes stratifying the random rankings so every node is favored, preventing zero estimated exposure in that construction. These statements concern the specified graph and randomization model. They do not validate an unknown bank interference graph or make reciprocal estimated probabilities unbiased. Probability approximation and its contribution to downstream error remain explicit qualification tasks.~\citep{exec25_cnet_paper,exec25_cnet_choice}

For 1-hop-max and bounded nonnegative responses, the selected text further bounds mean-outcome HT variance using $|B_4(i)|$ and $\pi_i^1$; under a restricted-growth condition it obtains a bound proportional to $n^{-1}(d_{\max}+1)^2\kappa^4/p$. Here $d_{\max}$ is maximum degree and $\kappa$ controls neighborhood growth. A bank cannot claim decreasing uncertainty merely by increasing its row count: degree, growth, exposure probabilities and outcome bounds matter. The text expressly distinguishes the available 1-hop-max variance proof from the randomized 3-net case. No unrestricted network-uncertainty guarantee is imported.~\citep{exec25_cnet_paper}

\subsection{An empty realized exposure arm is different from positive design support}

Positive $\pi_i^a$ for every borrower does not guarantee that both exposure arms contain observed borrowers in each experiment. This distinction matters for the ratio form of the H\'ajek estimator,
\[
 \widehat\mu_a^{\mathrm{H}}=
 \frac{\sum_i\mathbf{1}\{E_i^a\}Y_i(\mathbf a)/\pi_i^a}
      {\sum_i\mathbf{1}\{E_i^a\}/\pi_i^a}.
\]
The selected implementation initializes each arm mean to zero and updates it only when its denominator is positive. This defines a numerical extension for an empty arm. It is not the same expression as an everywhere-defined ratio, and zero is not an observed control outcome. The proposed result type should retain both exposed counts, both denominators and the convention or refusal applied.~\citep{exec25_cnet_rgcr,exec25_cnet_choice}

Consider three borrowers in one cluster, with potential outcomes $Y(\mathbf0)=(1,2,3)$ and $Y(\mathbf1)=(2,3,4)$ in a hypothetical monetary unit. The true average effect is one. Assigning the whole cluster with probability $1/2$ gives HT effect estimates $-4$ and $6$, each with probability $1/2$, so their mean is one. The source-style zero extension of H\'ajek gives $-2$ when everyone is controlled and $3$ when everyone is treated, yielding mean $1/2$. Every node has positive exposure probability in both arms, yet one realized arm is always empty. This exact example concerns the stated convention, not a native software run.~\citep{exec25_cnet_choice}

The same issue occurs in a small nontrivial randomized partition example. On the path $1-2-3$, enumerate the six strict orderings of independent continuous node scores and assign each node to its largest one-hop score. Across the induced partitions and independent half-probability cluster assignments, there are 20 weighted cases. The marginal full-exposure probabilities in each arm are $(5/12,1/3,5/12)$. With the preceding potential outcomes, exact enumeration gives HT mean one and variance $2242/75$; an independent joint-probability covariance expansion gives the same variance. The zero-extended H\'ajek mean is again $1/2$. These calculations both verify the intended exposure law and expose the effect of the empty-arm convention.~\citep{exec25_cnet_choice}

The selected 2020 text states a constant-effect unbiasedness result for H\'ajek and proves it by separating a constant treatment effect from each realized ratio. That algebra requires a nonzero realized denominator. The proof does not establish the source code's zero-extension guarantee on the assignments just enumerated. The proposal therefore retains this as a precise qualification boundary. Possible routes include a typed unavailable result, an explicitly redefined estimand or estimator, or an assignment design that guarantees the required denominators. Conditioning on successful experiments would itself change the sampling law and cannot be introduced without analysis. The later publisher version remains a separate source-review obligation.~\citep{exec25_cnet_paper,exec25_cnet_choice}

\subsection{RGCR integration requires more than recompilation}

The RGCR loader first looks for a binary graph cache and otherwise reads an edge list and writes that cache. It removes self-edges and requires consecutive node identifiers, connectivity and a loopless graph. These are explicit, useful restrictions, but they exclude some banking populations unless an adapter makes a justified transformation. A disconnected borrower network should not be made connected by inventing edges. A cached graph should not silently survive a changed source edge list. Stable borrower-to-node mappings, graph hashes and transformation receipts belong in the request.~\citep{exec25_cnet_rgcr,exec25_cnet_choice}

Randomness and input validation are additional gates. The inspected utility seeds the C generator from the clock; clustering and assignment use draws from that generator, while parallel simulation loops create work across threads. A safe analysis distinguishes reproducible seed input, draw allocation, thread scheduling and statistical independence. The build disables assertions, so an assertion in a parser or weight check is not a production rejection path. The weighted-score helper evaluates a logarithm of a draw that can be zero and divides by supplied weights; positive finite weights, endpoint handling and finite-score ties require explicit tests. A mixing-count helper also accesses the last element of a vector that can be empty for requests below ten. A proposed wrapper should reject unsupported counts before entering that path.~\citep{exec25_cnet_rgcr}

The simulation compares complementary assignments using the same partition. That pairing can reduce or change variability, but it means the two outputs are not automatically independent observations. The reported standard-error arithmetic needs a comparison against replicate-level or pair-level uncertainty on frozen small problems. Analytic variance routines also use supplied potential responses and joint exposure probabilities: this is valuable for known-law public evaluation, not an immediately available variance oracle for observed bank outcomes.~\citep{exec25_cnet_rgcr,exec25_cnet_choice}

The proposed decision is conditional algorithmic reference and controlled integration after these contracts are tested. No license-named file appears in the complete selected RGCR tree, so public availability is not treated as a redistribution grant. The old SNAP dependency, full artifacts and rights must be resolved before adoption. An independent reimplementation would need its own rights review and semantic tests; changing language does not remove the assignment, exposure or statistical obligations.~\citep{exec25_cnet_choice}

\subsection{NetworKit and SNAP: preserve identities, multiplicity and weights}

NetworKit's selected graph API supports optional weights and directed graphs but warns that multiple edges are not handled consistently. Its duplicate check is optional and disabled by default. The implementation confirms that enabling the check rejects a repeated edge; otherwise it appends adjacency entries. A borrower graph derived from repeated transactions must therefore choose an aggregation rule before admission. Counting transactions, summing monetary weights and deduplicating relationships define different exposure operators. The adapter should never rely on an unsupported duplicate representation merely because insertion succeeds.~\citep{exec25_cnet_networkit}

Deleting a NetworKit node leaves a hole in its identifier space; the actual node count and the upper identifier bound are separate values. SNAP likewise maintains a maximum assigned identifier bound independently of the number of remaining nodes. In an authored example with surviving IDs zero and two, there are two units, not three. If each has outcome one, dividing a zero-filled three-slot vector by its length produces $2/3$ instead of one. Compacting is a valid representation change only when outcomes, assignments, graph structure and returned results all use the saved forward and inverse maps.~\citep{exec25_cnet_networkit,exec25_cnet_snap,exec25_cnet_choice}

NetworKit's selected induced-subgraph and remapping constructions carry edge weights explicitly. Its restoration helper, however, traverses neighbors and calls edge insertion without the weight argument, whose default is one. Thus a node-ID restoration is not a guaranteed weighted-graph round trip in this selected path. If two incoming relationships have weights nine and one and neighbor treatments one and zero, the normalized weighted exposure is $9/10$; replacing both weights by one changes it to $1/2$. A public round-trip test should check node IDs, edge direction, multiplicity, weights and attributes separately. The finding is tied to this helper and revision, not asserted for every graph operation.~\citep{exec25_cnet_networkit,exec25_cnet_choice}

SNAP offers simple undirected and directed graph types and a separate directed multigraph. In the selected simple graph, duplicate insertion returns $-2$, while a successful insertion without an edge ID returns $-1$. Treating every negative return as failure would therefore misclassify successful operations. Its breadth-first traversal takes explicit incoming and outgoing direction flags and increments hops. The proposed wrapper should translate these API-specific conventions into typed outcomes and a declared distance definition. A weighted credit-exposure distance cannot be inferred from a hop traversal.~\citep{exec25_cnet_snap}

The libraries also expose different random-state contracts. SNAP's selected random generator treats zero as a request for timing-derived state; a positive supplied seed follows a separate path. NetworKit uses a thread-local engine seeded either identically across threads or with the thread identifier added, according to its explicit option. Equal seeds and equal-position draws can produce copied streams, while changing thread count or work scheduling changes how draws are assigned. In a hypothetical two-draw Bernoulli experiment, the variance of their mean is $1/8$ for independent draws but $1/4$ for identical draws. Thread safety alone does not settle this scientific distinction.~\citep{exec25_cnet_snap,exec25_cnet_networkit,exec25_cnet_choice}

Both remain conditional direct-integration candidates for specified graph services. NetworKit's MIT declaration and SNAP's BSD-style conditions support concrete artifact review, while dependency rights and installed behavior remain open. A port should be justified by a required runtime or representation constraint and tested against the same graph receipts. There is no reason to rewrite these graph algorithms merely to claim an internally owned causal library.~\citep{exec25_cnet_networkit,exec25_cnet_snap,exec25_cnet_choice}

\subsection{Repast HPC: distributed copies are not independent borrowers}

Repast HPC represents a shared network through local agents, nonlocal copies and complementary edges across processes. Its agent identity consists of original identifier, starting process and type; the current process is excluded from equality. That supports migration without changing identity, but it also means a bank evaluation must count scientific units independently of process-local copies. The selected graph stores one edge per ordered endpoint pair; a duplicate replaces the edge when allowed or marks the existing edge as conflicted. This behavior is different from a transaction multigraph and needs an explicit receipt.~\citep{exec25_cnet_repast}

The directed-neighbor delegate appends incoming and outgoing neighbors into the output vector. A reciprocal neighbor can appear twice, because the shared item helper does not deduplicate. This differs from the set-based generic adjacency inspected in Repast Simphony. Porting a model between the two frameworks therefore requires a semantic comparison rather than matching method names. Directional, unique-neighbor and incident-edge exposure should be separate adapter operations.~\citep{exec25_cnet_repast,exec25_cnet_choice}

The schedule comparator orders events by tick only, and the execution loop invokes same-tick callbacks sequentially. There is no explicit secondary phase key in that comparator. A repeating event with zero interval can re-enter the same tick indefinitely under the inspected loop, so positive finite intervals and future-time progress should be enforced at the boundary. Across processes, the runner selects the global minimum next tick. That synchronization does not itself refresh agent attributes or network projections. The source provides separate synchronization operations and expressly notes that synchronizing agent state does not update projection information.~\citep{exec25_cnet_repast}

Seed initialization is equally specific. A global-seed property can intentionally give processes the same seed, while the local-seed path derives rank-dependent state from the configured seed and records the effective value. These options serve different coupling purposes. Migration, rank count and assignment of draws to agents must be recorded if a public experiment is expected to replay. Agent ownership, graph-projection freshness and random-state ownership together define the simulated transition law. A valid MPI exchange is not evidence of causal or statistical independence.~\citep{exec25_cnet_repast,exec25_cnet_choice}

The decision is conditional use for distributed known-law experiments when scale justifies MPI, serialization and synchronization costs. Its BSD-style core terms do not clear every Boost, NetCDF or other installed dependency. Qualification should compare one-process and multi-process scientific outputs under a deliberately fixed coupling, exercise migration and stale-projection faults, and record refused or stalled requests. Small graph enumeration should precede expensive distributed runs.~\citep{exec25_cnet_repast,exec25_cnet_choice}

\subsection{igraph's C interface is a distinct interoperability option}

igraph is included explicitly as a C library usable from a C++ system. The selected neighbor API exposes direction, loop treatment and parallel-edge multiplicity, including a choice to collapse multiple edges. These are positive capabilities for a causal exposure adapter. The exact signature belongs to the pinned revision; a wrapper must not assume that an older installed version has the same interface. Vertex deletion reindexes IDs and invalidates iterators, and the induced-subgraph APIs expose mappings. Borrower identity should therefore live outside transient integer positions and be restored through the recorded maps.~\citep{exec25_cnet_igraph}

Its selected Dijkstra implementation validates weights, rejects negative values and NaNs, and documents explicit unreachable distances. That is useful failure behavior, but a caller must still decide whether a monetary relationship weight is a distance, a strength or an inverse cost. The proposed request includes units and transformation rules. A strong exposure relationship is not automatically a long path, and changing this interpretation can reverse which borrowers are considered close.~\citep{exec25_cnet_igraph,exec25_cnet_choice}

Random ownership requires attention even where a variable is thread-local. The selected default-pointer declaration is thread-local, but it initially points to a shared default generator object and backing state. Setting a custom default changes the pointer and requires that object to remain alive. A caller cannot infer private generator ownership merely from the pointer's storage class. The proposed native fixture should install and retain task-owned generators, check concurrent use and restore prior defaults deliberately. This is a source-based ownership finding, not a reported concurrency benchmark.~\citep{exec25_cnet_igraph}

The C interface can be adopted directly where its representation and failure contracts fit. Its GPL 2-or-later source declaration requires a distribution review for the intended product; porting or placing a process boundary around it is not declared an automatic rights solution. Complete intended artifacts, dependencies and native interface behavior remain qualification gates.~\citep{exec25_cnet_igraph,exec25_cnet_choice}

\subsection{A common public experiment contract}

The CLI and MCP surface should bind the scientific target, exposure model, graph and node-map hashes, edge units and multiplicity, event and availability clocks, partition generator, realized partition, conditional or marginal probabilities, random streams, worker layout and numerical tolerances. The response should preserve exposed-arm counts, probability uncertainty, graph transformations, typed failures, unsupported targets and complete cost denominators. An LLM may help propose this specification, but an unreviewed tool choice must not convert a graph service into a causal estimator or silently choose the empty-arm convention.~\citep{exec25_cnet_choice}

The public qualification sequence is concrete. First enumerate tiny graphs and all assignment cases, comparing marginal and joint exposure probabilities, HT moments and ratio failure behavior. Next test representation changes: duplicated transactions, removed nodes, compacted identifiers, weighted restoration and directed paths. Then freeze random-state and scheduling rules across threads or ranks, including stale graph caches and stale distributed projections. Finally compare matched native implementations with all attempts, failures, refusals and accepted outputs recorded. These tasks can proceed independently of unavailable bank truth, while expensive bank trials remain conditional on demonstrated public merit.~\citep{exec25_cnet_choice}

With a hypothetical USD 30,000 study budget, 1,000 attempts, 600 completions and 400 accepted outputs imply costs of USD 30 per attempt, USD 50 per completion and USD 75 per accepted output. None is a measured return or a price quote. They illustrate why a fast executable with frequent invalid exposure results can be economically inferior to a slower method with an appropriate scientific contract. The known repair routes are exact-law testing, explicit graph and random-state custody, typed refusal, rights review and matched native evaluation. Unknown bank interference, unsupported assignment truth and unrestricted dependent-data inference still require information or justified restrictions; this review does not establish complete C++ network qualification.~\citep{exec25_cnet_choice}


% SOURCE docs/release/2026-10-29/source/docs/writing-repair/2026-10-29/src/9a61f7787de4fb1a.tex CHAR 0
\section[Shared-state and remaining-language networks]{Shared-state inference and network contracts in Julia, Rust and JavaScript}
\label{sec:exec25-onet}

A bank sharing servicing capacity across arriving borrowers faces a different
interference problem from a bank observing a fixed borrower graph. One offer
may change the queue, available specialist time or a public price, affecting
later borrowers. A graph representation does not identify this mechanism, and
an accurate queue simulator does not establish the effect of changing policy.
Julia supplies direct research code for shared-state inference and components
for graph and dynamical simulation. Rust supplies graph representations with
explicit mutation contracts. JavaScript can carry a graph into an analyst
interface, but object and serialization semantics must preserve scientific
identity. We compare these roles before specifying public qualification.

The selected originals are pinned before the 7 October 2026 Shanghai cutoff:
Graphs.jl, Agents.jl, NetworkDynamics.jl, petgraph, Graphology and dml4ssi.
A second direct search lead is screened as a cost-to-equilibrium replication
package rather than accepted merely because its name contains spillovers.
Five successful bounded repository searches return six unique leads; three
additional searches fail with rate limits. Those failures leave coverage open;
they do not establish absence of Rust or JavaScript causal estimators. The
computations below are authored finite examples, not executions of acquired
packages.~\citep{exec25_onet_choice}

\subsection{Immediate treatment and changed shared state are different targets}

Write $X_t$ for pre-treatment borrower characteristics, $D_t\in\{0,1\}$ for a
servicing offer, $H_t$ for the shared state available before that offer, and
$Y_t$ for an outcome on a declared monetary and temporal scale. A direct effect
changes $D_t$ while retaining others' assignment law and the natural
shared-state distribution. A global effect changes the assignment history.
Hays and Raghavan's selected 2025 author version formalizes these distinct
targets. The relevant dependence must be mediated by the state. Their general
theory requires a homogeneous Markov structure with geometric ergodicity and
detailed balance, or a known finite dependence bound, alongside orthogonality,
moment and nuisance-rate conditions. A queue measured after the offer, an
omitted common liquidity shock or unrestricted persistence does not acquire
these properties by being entered as a feature.~\citep{exec25_onet_paper}

For the average direct effect, let $f(d,x,h)$ estimate the conditional outcome
and $e(x)$ the propensity in the selected model. The score is
\[
 \phi_t=f(1,X_t,H_t)-f(0,X_t,H_t)
 +\left\{\frac{D_t}{e(X_t)}-\frac{1-D_t}{1-e(X_t)}\right\}
   \{Y_t-f(D_t,X_t,H_t)\}.
\]
\[
 \widehat\tau=T^{-1}\sum_{t=1}^T\phi_t.
\]
The state is held fixed in each counterfactual prediction but marginalized
over in the target. This is not the effect of treating every borrower when
that policy changes capacity. The paper requires true and estimated assignment
probabilities bounded away from zero and sufficiently small nuisance errors.
Clipping a fitted propensity does not prove overlap of the bank mechanism.
\citep{exec25_onet_paper}

The paper uses independent auxiliary training data, explaining why ordinary
folds from a dependent sequence need not provide independence. Independently
generated trajectories offer a transparent first public experiment. A later
blocked temporal split needs its own dependence argument and cost accounting.
Declare the auxiliary population, outcome clock and measured state together;
neither a flexible forest nor prediction accuracy supplies causal restrictions.
\citep{exec25_onet_paper,exec25_onet_choice}

For switchback memory $m$, the global score reconstructs states under all-one
and all-zero histories and weights residuals by treatment-path probabilities.
Equation 5.7 of the primary text uses $D_{t-m:t}$: the preceding $m$ treatments
and the current treatment. Path positivity and state reconstruction are parts
of identification. An individual propensity or an observed-state prediction
cannot silently replace them. The theorem is conditional and asymptotic; it
does not certify every switching calendar or general network.
\citep{exec25_onet_paper}

\subsection{Reconcile the direct implementation with its equations}

dml4ssi exposes the simulated law, scores and comparisons, making it a useful
direct candidate. Its manifest identifies Julia 1.11.4. The selected nine-file
tree has no license-named original; reuse rights and complete installed
dependencies remain unresolved. Figure scripts establish globals and include
the inspected delegates. Some scalar convenience wrappers have signature or
return-shape concerns, but active simulation paths call vector-score helpers.
An unused wrapper concern must not be presented as evidence that every figure
fails.~\citep{exec25_onet_dml}

Several active expressions need reconciliation before adoption. The IPW
comparator averages weighted outcomes within each realized arm, whereas the
paper's Horvitz--Thompson comparator averages the signed score over all $T$
observations. For $(D,Y)=((1,2),(0,1))$ and $e=1/2$, these expressions return
$2$ and $1$. This exact witness establishes a normalization difference for
the selected input, not a package-wide failure rate or a refutation of the
paper. Preserve both expressions under distinct names and determine the
intended comparison.~\citep{exec25_onet_dml,exec25_onet_paper,exec25_onet_choice}

The switchback source checks $D_{\max(t-m+1,1):t}$ and uses
\[
 q_{\rm code}=\frac12-\frac{m-1}{4\ell},
\]
where $\ell$ is block length. Away from the initial boundary this window has
$m$ observations; the paper's displayed history has $m+1$. In an authored
stationary uniform-phase design with independent fair block assignments, a
window of length $L\leq\ell+1$ lies within one block for $\ell-(L-1)$ phases
and crosses a boundary for $L-1$ phases. Its all-one probability is
\[
 q_L=\frac{\ell-(L-1)}{\ell}\frac12+
       \frac{L-1}{\ell}\frac14
     =\frac12-\frac{L-1}{4\ell}.
\]
For $m=2,\ell=4$, enumeration gives $7/16$ for the source's two-observation
event and $3/8$ for the paper's three-observation history. The first truncated
single-observation event has probability $1/2$. The fixed expression is not
a universal initial-boundary correction. With $m=10,\ell=4$ it gives
$-1/16$, although enumeration of the ten-observation event gives $7/64$.
That last case is outside the paper's sufficient switching construction; it
motivates validation and a design-probability calculator, not a theorem
extension.~\citep{exec25_onet_dml,exec25_onet_paper,exec25_onet_choice}

A dependence horizon also needs one meaning in the state, assignment and
variance. The source's finite-lag loop uses lags strictly below its parameter
$m$; the paper's displayed variance includes lag $m$. For authored scores
$(1,-1,1,-1)$ and $m=2$, these long-run variance expressions are $-1/2$ and
$1/2$. The simulation clamps a negative mean-variance expression to zero
before taking its square root, producing a zero-width interval for this
particular witness. Clamping does not establish calibration. These calculations
compare finite expressions under declared indexing; a prototype must first
justify its actual dependence horizon, including assignment blocks, rather
than merely incrementing a loop bound.
\citep{exec25_onet_dml,exec25_onet_paper,exec25_onet_choice}

The Markov branch fixes square-root block counts. The paper requires a block
exponent exceeding $(1+\delta/2)^{-1}$ for its moment parameter $\delta$.
An exponent of $1/2$ needs a compatible stronger moment choice; $\delta>0$
alone is insufficient. Finite-size qualification must also require enough
complete blocks and state how unused tail observations are handled. A selected
public law can meet such conditions without implying that every persistent
bank sequence does.~\citep{exec25_onet_dml,exec25_onet_paper}

The IPW variance helper already divides by $T$, and the active caller divides
again. For four centered scores $(-2,2,2,-2)$, once- and twice-scaled variances
are $1$ and $1/4$, with standard errors $1$ and $1/2$. One ADE comparator
selects its first score before averaging; $(1,3,5)$ distinguishes that value
from the full mean. One switchback comparator gives a within-trajectory
variance helper an across-replication array. The DGP divides noise by
\texttt{sd\_error}, so an argument of two means standard deviation $1/2$
in that expression. Global-arm reconstructions draw separate initial states,
while observed trajectories follow their own initialization. Freeze the full
experiment definition, truth estimation, random streams and failure accounting
before using its comparison as investment evidence.
\citep{exec25_onet_dml,exec25_onet_choice}

Retain the method as a conditional reference, resolve rights and implement a
transparent public fixture with independent score, probability and covariance
oracles. Compare a pinned native Julia run only after those gates. Port if a
supported deployment language or measured full-workflow benefit requires it;
changing the interface does not resolve the conventions. No acquired code has
been executed here.~\citep{exec25_onet_choice}

\subsection{Graph positions and simulated agents are different scientific units}

Graphs.jl's selected simple graphs reject duplicate edge insertion but permit
self-loops. Adjacency access may return internal lists; copying duplicates
adjacency storage. Vertex removal swaps the last vertex into the removed
position, while induced subgraphs return new-to-old maps. Removing $b$ from
borrowers $(a,b,c)$ leaves positions $(a,c)$; a stale treatment array can
attribute $b$'s treatment to $c$. Persistent keys and explicit maps must travel
with transformations. The selected BFS takes the incoming-neighbor branch
for any direction token other than \texttt{:out}; validate direction tokens.
Weighted distances use a separate distance matrix that must follow the same
mapping.~\citep{exec25_onet_graphs,exec25_onet_choice}

Agents.jl's GraphSpace puts potentially many agents at a graph position. Its
documentation distinguishes that model from one-node-per-agent social
networks. A position query includes co-located agents; its agent-specific form
removes the focal agent. Two borrowers at one branch are therefore not two
adjacent borrower nodes. Choose the scientific unit before defining an exposure
denominator. Vertex removal deletes agents at that position and swaps stored
position lists. The selected body does not explicitly rewrite moved agents'
position fields. A native invariant check must compare membership, stored
location and identity after mutation; this is a scoped static observation,
not an observed runtime failure.~\citep{exec25_onet_agents,exec25_onet_choice}

StandardABM updates scheduled agents sequentially and places the model step
before or after them by configuration. Fastest, ID-sorted and randomized
orders are distinct. For two initially zero states and the authored rule
$x_i\leftarrow x_j+1$, forward order gives $(1,2)$, reverse order gives
$(2,1)$ and simultaneous old-state updates give $(1,1)$. A port that changes
this order changes the experiment. Bind schedule, RNG ownership and read/write
timing as scientific inputs.~\citep{exec25_onet_agents,exec25_onet_choice}

The selected checkpoint fields contain agents, space, properties, RNG and
maximum ID, but not stepping functions or elapsed model time. Deserialization
constructs StandardABM with a default scheduler unless supplied. Documentation
also explicitly excludes functions. RNG restoration alone cannot certify
experimental restart. Bind code, scheduler, clock, model class, outstanding
events and external data, then compare interrupted and uninterrupted
trajectories. Agents.jl is a conditional public simulation substrate here,
not an admitted general recovery or causal capability.
\citep{exec25_onet_agents,exec25_onet_choice}

\subsection{Network dynamics needs topology and state ownership}

NetworkDynamics.jl distinguishes simple graphs from ComponentGraph, which
preserves directed parallel edges and supplied order. Degree counts edges
with multiplicity; neighbor lists contain unique vertices. Edge-count and
unique-borrower exposures differ even on the same object. Component order
must match graph iterators; graphless construction can reorder legacy
components or select ComponentGraph when multiplicity requires it. Retain
edge identity, endpoints, units, component function and aggregation rule.
\citep{exec25_onet_dynamics}

Repeated component objects alias by default unless de-aliasing is requested.
Component copies de-alias metadata but share other fields; network copies
share topology. Captured mutable function state does not become independent
counterfactual state through that copy. Comparing capacity policies needs
explicit parameter, random-input and buffer ownership. Buffered, threaded and
accelerator execution, and sequential or threaded reducers, describe numerical
mechanisms; they do not establish the bank law, tolerance or uncertainty.
\citep{exec25_onet_dynamics,exec25_onet_choice}

Direct Julia integration is a conditional route for a prototype using explicit
dynamical components. A process interface can instead exchange immutable
topology, state and results. Both need small exact network oracles, units and
sign checks, solver and event settings, state-isolation tests and full costs.
Compare parallel execution only after these semantics agree.
\citep{exec25_onet_choice}

\subsection{Rust representations can preserve or collapse exposure}

petgraph Graph uses compact indices and moves the last node on removal.
StableGraph preserves unrelated indices but can reuse a freed slot. Neither
makes an integer a permanent borrower identity. Graph allows parallel edges;
GraphMap identifies nodes by value and updates an existing endpoint pair.
Conversion to GraphMap explicitly merges equal node weights and retains only
the last parallel edge. Use it only under established uniqueness and
multiplicity conditions.~\citep{exec25_onet_petgraph}

Suppose $a$ has two channels to treated $b$ with weights nine and one and a
channel of weight one to untreated $c$. Channel-weighted exposure is $10/11$.
Keeping only the last $a$--$b$ edge changes it to $1/2$. Summing parallel
weights could preserve this exposure but still lose channel attributes.
Summation, last-edge retention and duplicate refusal are distinct product
decisions. A port must not choose implicitly.
\citep{exec25_onet_petgraph,exec25_onet_choice}

The workspace declares Rust edition 2024, minimum Rust 1.91 and
MIT-or-Apache-2.0 licensing. These source declarations do not qualify an
installed binary or all dependencies. Dijkstra requires nonnegative costs
from a caller closure; valid exposure weights are not necessarily distances.
A conditional backend should expose typed failure, persistent external keys,
multiedge policy and validated costs. Prefer direct use or a narrow interface
when justified; any port still needs semantic and artifact qualification.
\citep{exec25_onet_petgraph,exec25_onet_choice}

\subsection{A JavaScript interface must preserve keys and immutable evidence}

Graphology's selected default is mixed, permits loops and disallows parallel
edges. Keys are coerced to strings: a numeric identifier and its string
rendering are not independent keys. Node data retains the supplied attributes
object. Export shallow-copies attributes, so nested objects remain shared.
In an authored example, changing an exported nested limit from 100 to 200 also
changes the original nested object. Export is not automatically an immutable
audit snapshot.~\citep{exec25_onet_graphology,exec25_onet_choice}

Use canonical namespaced keys, validated attributes, immutable serialized
snapshots and explicit graph options on every roundtrip. A displayed edge
must retain direction, type, evidence, availability clock and whether it is
a causal assertion or an observed interaction. An LLM-proposed edge belongs
in a separate editable layer until reviewed. A CLI or MCP request should bind
snapshot hash, target, exposure map, assignment law and method version, then
return a typed result or refusal. This comparison specifies those contracts;
it does not establish an implemented interface.
\citep{exec25_onet_graphology,exec25_onet_choice}

\subsection{One public qualification sequence connects the components}

Start with a small known-law servicing system and explicit state and assignment
calendar. Independent auxiliary trajectories train nuisances. Exact enumeration
checks treatment paths and support before larger simulations measure bias,
coverage, precision and refusal. A second fixture represents identical borrower
interactions in each graph backend and checks keys, direction, weights,
multiplicity and attributes through mutation and serialization. A third checks
sequential updates, counterfactual isolation and restart. Each answers a
different question; numeric agreement on one fixture cannot certify the others.
\citep{exec25_onet_choice}

Count all attempted, failed, refused and admissible requests. At an explicitly
hypothetical USD 30,000 cost, 1,000 attempts, 600 returned results and 400
admissible results imply USD 30 per attempt, USD 50 per returned result and
USD 75 per admissible result. Include acquisition,
conversion, auxiliary training, initialization, debugging, review and recovery.
Public/generated evidence can justify later bank experiments without supplying
authentic assignment truth or bank economics.~\citep{exec25_onet_choice}

Retain dml4ssi as a conditional method reference pending equation reconciliation,
rights and native tests. Consider Graphs.jl and petgraph as substrates,
Agents.jl for explicitly scheduled agents, NetworkDynamics.jl for declared
dynamics and Graphology for a reviewed interface. Reject silent graph collapse,
unverified state copying and automatic causal admission. Bounded search tails and failed queries remain separate source-coverage
obligations. Unknown state
sufficiency, hidden confounding and general dependent-data uncertainty require
information or justified restrictions. These originals do not establish
complete remaining-language network qualification.~\citep{exec25_onet_choice}

\subsection{The allocation law determines the comparison}

The chapter's alternatives cannot be ranked by the presence of a graph object
or a spillover label. A spatial equilibrium, a randomized exposure contrast,
a grouped policy, a learned allocation and a shared-state intervention use
different information. Their comparison must retain the population, assignment
law, exposure support, background units and uncertainty target. Graph identity,
update order and random-state ownership matter because they can change those
scientific objects during computation.

The practical choice is conditional reuse of a component whose target and
failure behavior match the banking request, with exact small examples supplying
a first comparison and native trials supplying separate implementation evidence.
The next chapter returns to customer histories, heterogeneity and monetary units,
where the same discipline is needed even when an explicit network is absent.


